% ----------------------------------------------------------------------
% Title:    aTableau.  Author:   Andrew Mathas Created:  October 2023
% Modified: 16:58 Friday, 24 January 2025
% ----------------------------------------------------------------------
\documentclass{amsart}
\usepackage[svgnames,table]{xcolor}
\usepackage{atableau}
\usepackage{xspace}
\usepackage{manfnt}
\usepackage{booktabs}
\usepackage{longtable}
\usepackage{xltabular}
\usepackage{siunitx}
\usepackage{array}
\newcolumntype{L}{>{\sffamily\color{MidnightBlue}\textbackslash}l}

% ----------------------------------------------------------------------
\usepackage{bm}
\usepackage{mathtools}

% ----------------------------------------------------------------------
\usepackage{bbding}
\newcommand\No{\textcolor{red}{\XSolidBrush}}
% \Yes, which adds a hyperlink, is defined below

% ----------------------------------------------------------------------
\usepackage[a4paper,hmargin=24mm, vmargin=18mm, footskip=10mm]{geometry}
\synctex=1
\hfuzz=5pt
\parindent=0pt
\parskip=4mm

\makeatletter
\def\@evenhead{\hfill\aTableau\hfill}
\let\@oddhead\@evenhead
\def\@oddfoot {\color{aTableauInner} Version \atableau@version~(\atableau@release)\hfill\thepage}
\def\@evenfoot{\color{aTableauInner} \thepage\hfill Version \atableau@version~(\atableau@release)}
\makeatother

% ----------------------------------------------------------------------
\newcommand\aTableauColour[1]{%
  \tikz[baseline={(current bounding box.south)}]\node[minimum width=8mm, minimum height=2ex,fill=aTableau#1]{};%
}

\ExplSyntaxOn
% \optionlabel applies \option to the csv in #1 except in the case
% when #1 starts with a * when the rest of #1 is printed
\NewDocumentCommand\optionlabel{ m }
{
    \str_if_eq:eeTF { \str_head:n {#1} } {*}
    {
        \tl_set:Nx \l_tmpb_tl { \str_tail:n {#1} }
        \tl_put_left:Nn \l_tmpb_tl {\sffamily\color{aTableauInner}}
    }
    {
        \tl_clear:N \l_tmpa_tl
        \tl_clear:N \l_tmpb_tl
        \clist_map_inline:nn {#1}
        {
            \tl_put_right:Ne \l_tmpb_tl { \l_tmpa_tl \exp_not:N\option{##1}}
            \tl_set:Nn \l_tmpa_tl {,~}
        }
    }
    \l_tmpb_tl\hss
}

\NewDocumentCommand\Yes{o d()}
{
  \IfNoValueTF{#1}
  {
    \IfNoValueTF{#2}
    { \hyperref[\g_lastlink_tl]{\textcolor{ForestGreen}{\CheckmarkBold}} }
    { \hyperref[Key:#2]{\textcolor{ForestGreen}{\CheckmarkBold}} }
  }
  {
    \hyperref[\g_lastlink_tl-\str_lowercase:f{#1}]{\textcolor{ForestGreen}{\CheckmarkBold}}
  }
}

% Add hyperinks for (new displayed) options
\tl_new:N \g_lastlink_tl
\prop_new:N \l__atableau_options_prop
\NewDocumentCommand\AddKeyLink{m}
{
   \tl_gset:Nx \g_lastlink_tl {Key:\str_lowercase:f{#1}}
   \prop_if_in:NVF \l__atableau_options_prop \g_lastlink_tl
   {
      \phantomsection\exp_args:Nf\label{Key:\str_lowercase:f{#1}}% add a label for references
      \prop_put:NVn \l__atableau_options_prop \g_lastlink_tl {}
   }
}
% \option*[default value]{option name}[accepted values]<hyperlink context>(link name)
% The \option command is slightly overloaded, both in terms of arguments
% (there are five) and in terms of use.
% 1. If #2 is given then \option gives the syntax for the using this key:
%    - #1=* determines whether a \newline or \bigskip is inserted. The
%      *-version is used when a block of keys are listed together
%    - #2 gives the default options
%    - #3 is the key name and optionally may take the form key=value
%      to describe the allowed values. If the value is given then it is
%      printed in orange after the key inside \langle...\rangle. This is
%      intended for "generic values like 'value' or 'style'.
%      The key component of #2 is used to add a hyperlink to the current
%      location, which takes the form Key:<key>-<context>, where the
%      <context> is added only if #5 is supplied
%    - any accepted values in #3 are printed in the right-hand margin
% 2. When #2 is omitted, this an inline citation of the key. A
% hyperlink is added

% The *-version does not do a \bigskip before the option. A hyperlink is
% added for new displayed options. If #2 is blank then the *-version
% puts (no) before the option key. If #5 is given then it is added to
% to the hyperlink.
\NewDocumentCommand\option{ somo d<> d()}
{%
  \seq_set_split:Nnn \l_tmpa_seq {=} {#3}%
  \seq_pop_left:NN \l_tmpa_seq \l_tmpa_tl%
  \seq_if_empty:NTF \l_tmpa_seq
  { \tl_clear:N \l_tmpb_tl }
  { \tl_set:No \l_tmpb_tl { \seq_use:Nn \l_tmpa_seq {=} } }
  \IfNoValueTF{#2}
  {%
    \IfNoValueTF{#6}
    {
      \tl_gset:Nx \g_lastlink_tl {Key:\str_lowercase:f{\l_tmpa_tl}}
      \IfNoValueF{#5} { \tl_gput_right:No \g_lastlink_tl {\str_lowercase:f{-#5}} }
    }
    { \tl_gset:Ne \g_lastlink_tl {Key:#6} }
    \hyperref[\g_lastlink_tl]{\textsf{\color{aTableauInner}\IfBooleanT{#1}{(no)~}\l_tmpa_tl}}
    \index{\l_tmpa_tl}%
    \tl_if_empty:NF \l_tmpb_tl { \textsf{\,$=$\,}\textsf{\textcolor{orange}{\l_tmpb_tl}} }
  }
  {%
    % encourage page breaks *before* the option, rather than after it
    \IfBooleanTF{#1}{\newline}{\bigskip\pagebreak[0]}%
    % adjust #3 if it contains a *, when we add = <value>
    % add a label for references
    \IfNoValueTF{#5} { \AddKeyLink{\l_tmpa_tl} }{ \AddKeyLink{\l_tmpa_tl-#5} }
    \noindent\textsf
    {
        \textcolor{aTableauInner}{\l_tmpa_tl}
        \tl_if_empty:NF \l_tmpb_tl { \textcolor{orange}{$\,=\,\langle$\l_tmpb_tl$\rangle$} }
    }
    \tl_if_blank:nF {#2}
    { ~ \tl_if_eq:nnTF{#2} {-} {(default)} {(default:~\textcolor{DarkRed}{#2})} }%
    \hfill%
    \IfNoValueF{#4}{\textcolor{Grey}{[accepts:~#4]}}%
    \index{\l_tmpa_tl|underline}%
    \nopagebreak[4]%
  }%
}

% index command ignoring the backslash and capitalisation
\cs_new_protected:Npn \l_addto_index:nn #1#2 { \index{#2@\textbackslash #1} }
\cs_generate_variant:Nn \l_addto_index:nn {Vo}
\NewDocumentCommand\indexcmd{m}
{
    \tl_set:Nn \l_index_tl  {#1}
    \l_addto_index:Vo \l_index_tl { \str_lower_case:n\l_index_tl }
}

\clist_new:N \l_atableau_keys_clist
\clist_set:Nn \l_atableau_keys_clist
{
  Australian,
  English,
  French,
  Ukrainian,
  align,
  australian,
  border,
  box,
  boxes,
  no boxes,
  color,
  colour,
  cols,
  conjugate,
  left delimiter,
  right delimiter,
  delimiters,
  dotted,
  e,
  east,
  english,
  entries,
  font,
  french,
  halign,
  height,
  inner,
  label,
  left,
  math,
  nodes,
  north,
  ribbon,
  ribbons,
  right,
  rows,
  scale,
  separator,
  separators,
  shape,
  shifted,
  skew,
  snobs,
  snob style,
  ribbon style,
  path style,
  snob box,
  ribbon box,
  path box,
  snob box style,
  ribbon box style,
  path box style,
  valign,
  south,
  star style,
  abacus star style,
  tabloid,
}
\NewDocumentCommand\HighlightKeys{m}
{
        \clist_if_in:NnTF \l_atableau_keys_clist {#1}
        {
            \textcolor{orange}{#1}
        }
        {
            #1
        }
}
\ExplSyntaxOff

% ---------------------------------------------------------------------------
\usepackage{enumitem}
\setlist{itemsep=2mm, parsep=2mm}
\newlist{options}{description}{1}
% remove the colon from the description label by redefining \makelabel
\setlist[options]{
  before=\let\makelabel\optionlabel,
  labelwidth=\textwidth,
  nosep,
}

% ---------------------------------------------------------------------------
% the index

\usepackage{imakeidx}
\indexsetup{level=\section*, toclevel=section, noclearpage}
\makeindex[columns=3]
\renewcommand*{\seename}{see}

% ---------------------------------------------------------------------------
%\renewcommand\thesubsection{\thesection\Alph{subsection}}
\setcounter{secnumdepth}{3}
\setcounter{tocdepth}{3}
\renewcommand\thesubsection{\thesection\Alph{subsection}}
\renewcommand\thesubsubsection{\thesubsection.\arabic{subsubsection}}

\makeatletter
\def\l@subsection{\@tocline{2}{0pt}{2.5pc}{5pc}{}}
\def\l@subsubsection{\@tocline{2}{0pt}{3.5pc}{3pc}{}}
\makeatother

\newenvironment{warning}{\trivlist\item[\hskip\labelsep\textcolor{red}{\textdbend}]}{\endtrivlist}

% ---------------------------------------------------------------------------
\usepackage[skins]{tcolorbox}
\usepackage{hyperref}
\hypersetup{
  colorlinks=true,
  linkcolor=MediumBlue,
  citecolor=FireBrick,
}
\tcbuselibrary{listings, breakable}

\DeclareTotalTCBox\keyword{ O{} v }
{
  bottom=0pt,
  colback  = Ivory,
  colupper = aTableauInner,
  fontupper= \bfseries,
  fontlower= \bfseries,
  left     = 0mm,
  right    = 0mm,
  size     = tight,
  skin     = tile,
  top      = 0pt,
  nobeforeafter,
  on line,
  shrink tight,
  verbatim,
  #1}
{#2}

\newcommand\aTableau{\textcolor{aTableauMain}{\sffamily aTableau}\xspace}

% hyperref links to ctan
\newcommand\ctan[1]{\href{https://www.ctan.org/pkg/#1}{\texttt{#1}}}
\newcommand\TikZ{\href{https://www.ctan.org/pkg/tikz}{Ti\textit{k}Z}\xspace}

\lstdefinestyle{atableau}
{
  style=tcblatex,
  classoffset=0,
  texcsstyle=*\color{MediumBlue},%
  deletetexcs={begin, end},
  commentstyle=\color{Peru},
  moretexcs={,% aTableau package commands
    aTabset,
    Abacus,
    Diagram,
    Multidiagram,
    Multitableau,
    RibbonTableau,
    ShiftedDiagram,
    ShiftedTableau,
    SkewTableau,
    SkewDiagram,
    Tableau,
    Tabloid,
  },%
%  classoffset=1, % aTableau package options
%  keywordstyle=\color{aTableauInner},%
%  morekeywords={
%    Australian,
%    English,
%    French,
%    Ukrainian,
%    align,
%    australian,
%    border,
%    box,
%    boxes,
%    color,
%    colour,
%    cols,
%    conjugate,
%    delimiter,
%    delimiters,
%    dotted,
%    e,
%    east,
%    english,
%    entries,
%    finite,
%    font,
%    french,
%    halign,
%    height,
%    infinite,
%    inner,
%    label,
%    left,
%    math,
%    no,
%    nodes,
%    north,
%    only,
%    ribbon,
%    ribbons,
%    right,
%    rows,
%    scale,
%    separator,
%    separators,
%    shape,
%    shifted,
%    skew,
%    snobs,
%    valign,
%    south,
%%   star,
%%   style,
%    tabloid,
%    text,
%    text,
%    type,
%    ukrainian,
%    wall,
%    west,
%    width,
%    xoffset,
%    xscale,
%    yoffset,
%    yscale,
%  },
%  classoffset=2, % aTableau values
%  keywordstyle=\color{orange},%
%  morekeywords={
%    A, C, DD, AA,
%    beta,
%    bottom,
%    center,
%    centre,
%    contents,
%    false,
%    first,
%    hooks,
%    last,
%    left,
%    residues,
%    right,
%    row,
%    top,
%    true,
%  },
%  classoffset=3, % TikZ and LaTeX commands
%  keywordstyle=\color{teal},%
%  morekeywords={
%    draw,
%    foreach,
%    node,
%    small,
%    tiny,
%    tikz,
%    tikzpicture,
%    tikzset,
%  },
%  classoffset=4,% TikZ keys
%  keywordstyle=\color{red},
%  morekeywords={
%    circle,
%    fill,
%    overlay,
%    picture,
%    radius,
%    remember,
%    style,
%  },
  classoffset=5,   % commands that we want to hide :)
  keywordstyle=\color{LightSkyBlue!10},
  morekeywords={
    bigskip,
    medskip,
    newline,
    par,
    qquad,
    quad,
    smallskip,
  },
}

\newcounter{examples}
\numberwithin{examples}{subsection}
\newcommand\examplesautorefname{Example}

\NewTCBListing[use counter=examples, number within=subsection,%crefname={Example}{Examples}
]{example}{ !O{} }{%
  enhanced,
  skin=bicolor,
  colframe=RoyalBlue!70,
  colbacklower=OldLace,
  colback=LightSkyBlue!10,
  colbacktitle=RoyalBlue!70,
  attach boxed title to top right={xshift=0mm,yshift=-3.1mm},
  boxed title style={
    size=small,
    arc=1mm,
    bottom=-1mm,
    sharp corners=northwest,
    sharp corners=southeast,
  },
  title={\tiny\thetcbcounter},
  sharp corners=northeast,
  parbox=false,
  lefthand width=0.34\textwidth,
  listing style=atableau,
  parbox=false,
  parskip=4mm
  sidebyside,
  sidebyside align=top seam,
  sidebyside gap=2mm,
  text and listing,
  text outside listing,
  boxsep = 0pt,
  left=0mm,
  #1
}

\let\thetcolorboxnumber\relax

\DeclareTCBListing[use counter=examples, number within=subsection,%crefname={Example}{Examples}
]{Example}{ !O{} }{%
  enhanced,
  skin=bicolor,
  colframe=RoyalBlue!50,
  colbacklower=OldLace,
  colback=LightSkyBlue!20,
  colbacktitle=RoyalBlue!70,
  breakable,
  width=\textwidth,
  attach boxed title to top right={xshift=0mm,yshift=-3.1mm},
  boxed title style={
    size=small,
    arc=1mm,
    bottom=-1mm,
    sharp corners=northwest,
    sharp corners=southeast,
  },
  title={\tiny\thetcbcounter},
  sharp corners=northeast,
  parbox=false,
  listing style=atableau,
  text and listing,
  listing above text,
  %text outside listing,
  center lower,
  #1
}

% ---------------------------------------------------------------------------
\NewDocumentCommand\set{ O{} >{\SplitArgument{1}{|}}m }{#1\{\,\SetContents{#1}#2\,#1\}}
\NewDocumentCommand\SetContents{mmm}{\IfNoValueTF{#3}{#2}{#2\,#1|\,\mathopen{}#3}}

\newcommand\T{\mathsf{t}}
\newcommand\blam{{\bm{\lambda}}}

% ---------------------------------------------------------------------------
% manual title
\makeatletter
\author{Andrew Mathas}
\usepackage{tikz}
\usetikzlibrary{
  backgrounds,
  decorations.pathmorphing,
  matrix,
  patterns,
  shadows.blur,
  shapes.geometric,
}
\tikzset{
  shadowed/.style={
    blur shadow={shadow blur steps=5},
    bottom color=RoyalBlue,
    draw=NavyBlue!70,
    shade,
    font=\normalfont\Huge\bfseries\scshape,
    %rounded corners=8pt,
    top color=LightSkyBlue!30,
  },
  boxes/.style={draw=RoyalBlue,
    fill=AliceBlue,
    font=\sffamily\small,
    inner sep=5pt,
    rectangle,
    rounded corners=8pt,
    text=MidnightBlue,
  }
}

% define the heading as a command because it contains \atableau@version
\newcommand\aTableauHeader{%
  \begin{center}
    \begin{tikzpicture}
      \aTabset{border=false}
      \draw[shadowed, rounded corners](30mm,0) rectangle node[white]{ } (\paperwidth-30mm,22mm);
      \begin{scope}[shift={(2.6,-0.2)}]
        \aTabset{scale=0.5, ribbon style={fill=LightSkyBlue!10}, inner wall=MidnightBlue}
        \Diagram(4,2) [ribbons={11,12,13,21,23,31,32,33,41,43,51,53}]{}     % A
        \Diagram(5,2) [ribbons={11,12,13,22,32,42,52}]{}                    % T
        \Diagram(6,2) [ribbons={11,12,13,21,23,31,32,33,41,43,51,53}]{}     % A
        \Diagram(7,2) [ribbons={11,12,13,21,23,31,32,33,41,43,51,52,53}]{}  % B
        \Diagram(8,2) [ribbons={11,21,31,41,51,52,53}]{}                    % L
        \Diagram(9,2) [ribbons={11,12,13,21,31,32,33,41,51,52,53}]{}        % E
        \Diagram(10,2)[ribbons={11,12,13,21,23,31,32,33,41,43,51,53}]{}     % A
        \Diagram(11,2)[ribbons={11,13,21,23,31,33,41,43,51,52,53}]{}        % U
      \end{scope}
      \node[anchor=west,boxes] at (4cm,0cm) {Andrew Mathas};
      \node[anchor=east,boxes] at (\paperwidth-4cm,0) {Version \atableau@version};
    \end{tikzpicture}
  \end{center}
}

% ----------------------------------------------------------------------
\hypersetup{
  pdfcreator={ Generated by pdfLaTeX },
  pdfinfo={
    Author  ={ Andrew Mathas },
    Keywords={ Algebraic combinatorics, representation theory, abacus, Young diagram, tableau, symmetric groups  },
    License ={ LaTeX Project Public License v1.3c or later },
    Subject ={ LaTeXing tableaux and friends },
    Title   ={ ATableau - \atableau@version }
  },
  colorlinks=true,
  linkcolor=blue,
  urlcolor=MediumBlue,
}

\makeatother

% ----------------------------------------------------------------------
\numberwithin{equation}{subsection}
\def\equationautorefname~#1\null{(#1)\null}
\renewcommand\sectionautorefname{Chapter}
\renewcommand\subsectionautorefname{Section}
\def\subsubsectionautorefname~#1\null{\S#1\null}

\begin{document}

  % --------------------------------------------------------------
  % title page
  \begin{center}

    \aTableauHeader

    \begin{quote}
      A \LaTeX\ package for symmetric group combinatorics, including Young
      diagrams, tableaux, tabloids, skew tableaux, shifted tableaux, ribbon
      tableaux, multitableaux, abacuses.
    \end{quote}

    \begin{minipage}{0.75\textwidth}
      \tableofcontents
    \end{minipage}

    \aTabset{scale=0.6}
    \tikzset{
       B/.style={fill=blue,text=white},
       G/.style={fill=ForestGreen,text=yellow},
       R/.style={fill=red,text=white},
       Y/.style={fill=yellow,text=ForestGreen},
       b/.style={ball color=blue},
       r/.style={ball color=red},
       w/.style={ball color=white},
    }
    \begin{tikzpicture}
      \Tableau(0,0){[R]1[B]2,[G]3,[Y]4} \quad
      \Tableau(2.25,1)[australian]{[R]1[B]2,[G]3,[Y]4} \quad
      \Tableau(3.5,1.75)[skew={0,2}]{[Y]4[G]3[R]1,[B]2} \quad
      \Tableau(6,0)[ukrainian,skew={0,1^2}]{[B]2[R]1,[G]3,[Y]4} \quad
      \Tableau(7.5,-1)[french,skew={0,1^2}]{[B]2[R]1,[G]3,[Y]4} \quad
      \Tableau(6.0,-2.25)[ukrainian]{[R]1[G]3[Y]4,[B]2} \quad
      \Tableau(3.5,-3.0)[french]{[R]1[G]3[Y]4,[B]2} \quad
      \Tableau(2,-1)[australian]{[B]2,[R]1[G]3[Y]4} \quad
      \Abacus(4,0.5)[abacus ends=--,scale=1.4]{3}{[b]9,[w]8,[r]8,[b]6,[b]6,[r]2,[b]1}
    \end{tikzpicture}

  \end{center}

  % --------------------------------------------------------------
  \section{Introduction}

  The package provides commands for drawing common objects in algebraic
  combinatorics and representation theory, together with styling options
  and a key-value interface for additional customization.  Under the
  hood, everything is drawn using \TikZ (and \LaTeX3).  These commands
  can be used either as standalone commands or as part of a
  \keyword{tikzpicture} environment.  This package started as code for
  drawing the diagrams that I need in my own research. It is written as
  flexibly as possible in the hope that it will draw the diagrams that
  others need.

  For the impatient, here is a basic example that shows how to use the
  two main commands provided by the \aTableau package:

  \begin{example}
    \Abacus{4}{4^3,2,1^2,0} \qquad
    \Tableau{1239,456{10},78}
  \end{example}

  If you already know what tableaux and abacuses are, then you can
  probably work out the (basic) syntax of these two commands from this
  example and start using the package. If you do not know what tableaux
  and abacuses are, then you probably don't need this package! In any
  case, you will not find the precise mathematical definitions of these
  objects in this manual, so please consult textbooks like
  \cite{Mathas:ULect,Macdonald} if you need them. For more advanced
  usage, such as adding styling, I recommend skimming through the
  manual, and looking at the many examples.  \autoref{S:keys} gives
  quick summary of all \aTableau commands and their options.

  The \aTableau commands are intended to be easy to use and easy to
  customise. For example, partitions can be entered either as a (weakly
  decreasing list of non-negative) integers, or using exponential
  notation, or as a sequence of beta numbers (for the \keyword{\Abacus}
  command), and tableaux are entered as comma-separated words. The
  commands support the main conventions for drawing tableaux and
  diagrams.

  By way of example, the following table shows how to draw a
  Young diagram of shape $\lambda=(4,3^2,2)$ using the four tableaux
  conventions supported by \aTableau:

  \begin{center}
      \aTabset{align=centre}
      \begin{longtable}{lp{0.6\textwidth}l} \toprule % longtable allows footnotes
        \textcolor{aTableauMain}{Convention} &\textcolor{aTableauMain}{\aTableau command} & \textcolor{aTableauMain}{Picture} \\ \midrule
          English    & \keyword{\Tableau[english]{1234,567,89{10},{11}{12}}}\newline(the default)
                     & \Tableau{1234,567,89{10},{11}{12}}\\[10mm]
          French     & \keyword{\Tableau[french]{1234,567,89{10},{11}{12}}}
                     & \Tableau[french]{1234,567,89{10},{11}{12}}\\[10mm]
          Ukrainian
          \footnote{Some authors call this the Russian convention. The \option{russian}(ukrainian) key may also be used.}
                     & \keyword{\Tableau[ukrainian]{1234,567,89{10},{11}{12}}}
                     & \Tableau[ukrainian]{1234,567,89{10},{11}{12}}\\[10mm]
          Australian
          \footnote{I have not seen a name for this convention in the literature. Calling this the \textit{Australian convention} seems apt.}
                     & \keyword{\Tableau[australian]{1234,567,89{10},{11}{12}}}
                     & \Tableau[australian]{1234,567,89{10},{11}{12}} \\[10mm]
          \bottomrule
      \end{longtable}
      \index{english@English}
      \index{french@French}
      \index{ukrainian@Ukrainian}
      \index{australian@Australian}
  \end{center}

  This table shows how the commands provided by \aTableau work: there
  are basic commands for drawing pictures, and these pictures can be
  embellished using optional arguments, which use a key-value interface.

  The different conventions for diagrams, and tableaux, are listed in
  order of (my perception of) their popularity in the literature.  By
  default, the \textit{English} convention is used, so we can omit
  \emph{english}:
  \begin{example}[lefthand width=0.4\textwidth]
  \Tableau{1234,567,89{10},{11}{12}}
  \end{example}
  for the first tableau in the table. These four conventions are used
  in exactly the same way with the \keyword{\Diagram} command:
  \begin{example}[lefthand width=0.4\textwidth]
  \Diagram{4,3,3,1} \qquad
  \Diagram[ukrainian]{4,3^2,1}
  \end{example}
  This example shows that partitions can either be typed out in full, as
  a comma-separated list of positive integers, or using exponential
  notation, which is usually shorter and easier to read.

  To use the package, add the following line to your document preamble:

  \indexcmd{usepackage}
  \begin{quote}
    \keyword{\usepackage[options]{atableau}}
  \end{quote}

  Here, the \textit{options} is an optional comma-separated list of
  \aTableau options, or keys. We describe the possible options below.
  For example, to make the tableaux and Young diagrams commands use the
  \option{French} convention, by default, you could write:

  \begin{quote}
    \keyword{\usepackage[french]{atableau}}
  \end{quote}

  You can change the default options at any point in your document using
  the \keyword{\aTabset} command, and you can change them for any
  diagram.
  \indexcmd{aTabset}

  Each of the \aTableau commands is designed to be easy to use, easy to read,
  and easy to customise. Each command accepts an optional list
  of key-value options that modify and embellish the diagrams. In addition, optional
  styling can be applied to entries of a tableau and the beads in an
  abacus.  Each command can be used as a standalone object in a
  document, or they can be incorporated into a larger
  \keyword{tikzpicture} environment, making it possible to use
  these commands to create more complicated pictures.

  The general syntax of an \aTableau command is:
  \begin{quote}
    \keyword{\Command (x,y) [options] {...specifications...}}
  \end{quote}

  \index{Cartesian coordinates}
  The $(x,y)$-Cartesian coordinates%
  \footnote{Currently, only Cartesian coordinates are supported. For
  example, it is not possible to use named or polar coordinates.}
  are optional, being necessary if (and only if), the diagram is part of
  a~\keyword{tikzpicture} environment. See \autoref{SS:coordinates} and
  \autoref{SS:abcoordinates} for more detail about how to use these coordinates
  for tableaux and abacuses, respectively.

  The key-value \emph{options}, or \emph{keys}, control the style and
  appearance of the final diagram. The keys are optional and are
  described in detail below, with examples. This section emphasises only
  some of these options together with the fact that every key can be set
  ``globally'' (inside the current \LaTeX\ group), by using
  \keyword{\aTabset{...options...}}, or ``locally'', for the current
  picture. As discussed already, the default options for the entire
  document can also be set via
  \keyword{\usepackage[options]{atableau}}.\indexcmd{aTabset}

  Many of the \aTableau options are specific to the particular diagrams
  being drawn. The full list of the \aTableau commands and their options
  is given in \autoref{S:keys}, with the following sections giving more
  details and many examples. The keys that are common to all
  \aTableau commands are:

  \begin{options}
    \item[align]\index{align}
    Sets the baseline for vertically aligning diagrams. The options
    are \keyword{align=top}, \keyword{align=center}, and
    \keyword{align=bottom}, which aligns the baseline of the diagram
    with the top, centre or bottom of the surrounding objects,
    respectively. By default, all \aTableau pictures are centered.

    \item[math entries, text entries]\index{math entries}\index{text entries}
    Determines whether tableaux entries, and bead labels, are typeset as
    mathematics or text.

    \item[name]\index{name}
    Set the prefix for the \TikZ \emph{named anchors} for the boxes and
    beads in \aTableau diagrams.

    \item[scale, xscale, yscale, script, scriptscript]
    \index{scale}\index{xscale}\index{yscale}\index{script}\index{scriptscript}
    Use \option{scale} to rescale all \aTableau diagrams. Use
    \option{xscale} to rescale in the $x$-direction and \option{yscale}
    to rescale in the $y$-direction. The \option{script} and
    \option{scriptscript} keys can be used to fine-tune the size of
    \aTableau commands when used as subscripts (rare) and subsubscripts
    (unlikely).

    \item[styles]\index{styles}
    A short-hand for defining single-use \TikZ-styles for the current
    picture.

    \item[tikzpicture]\index{tikzpicture}
    When used without Cartesian coordinates, all \aTableau commands are
    implicitly drawn inside a \keyword{tikzpicture} environment.  Use
    \option{tikzpicture=keys} to set the optional argument of this
    environment: \keyword{\begin{tikzpicture}[ keys ]...\end{tikzpicture}}.

    \item[tikz before, tikz after]\index{tikz before}\index{tikz after}
    Use the \option{tikz before} and \option{tikz after} keys to inject
    \TikZ code into the \keyword{tikzpicture} environment, before and
    after the picture drawn by \aTableau.

  \end{options}

  With boolean options, which are set to \textsf{true} or
  \textsf{false}, the value can typically be omitted. For example,
  \option{border} is equivalent to \option{border=true} and to
  \option{no border=false}. The \option{no border} key is inverse to
  \option{border}, so \option{no border} is equivalent to \option{no
  border=true}, and to \option{border=false}.

  \index{American dialect}
  For all keys, American spelling variations are tolerated. For example, you
  can use \emph{center} and \emph{color} instead of the correct
  spellings of \emph{centre} and \emph{colour}, respectively.

  \begin{warning}
    Unlike the egalitarian \aTableau settings, \TikZ-style settings
    \emph{do not} countenance English spelling of the English language!
  \end{warning}

  The next example contrasts using the \keyword{\Diagram} command both
  outside and inside a \keyword{tikzpicture} environment. More
  explicitly, the second \keyword{\Diagram} command has Cartesian
  coordinates because it is being used inside a \keyword{tikzpicture}
  environment. Otherwise, these two \keyword{\Diagram} commands are
  identical, except that the second diagram is embellished with labels
  that are drawn by the other commands in the \keyword{tikzpicture}
  environment.

  \begin{Example}
  % \usetikzlibrary{patterns}
  \tikzset{
    B/.style={pattern=north east lines, pattern color=blue},
    G/.style={pattern=north west lines, pattern color=ForestGreen},
  }
  \Diagram[ribbons={(G)16ccccrrr, (B)16crrcccr}]{6,5^2,2} % outside tikzpicture
  \qquad \qquad
  \begin{tikzpicture} % inside tikzpicture with labels
    \Diagram(0,0)[ribbons={(G)16ccccrrr, (B)16crrcccr}]{6,5^2,2}
    % add the labels
    \draw[->](1.5, 0.5)node[right]{$\alpha=(1,2)$} to [out=180,in=90]  (A-1-2.base);
    \draw[->](1.5,-2.5)node[right]{$f_\alpha=(4,2)$} to [out=180,in=270] (A-4-2.base);
    \draw[red,<->]([xshift=-1mm]A-1-1.south west)
       --node[rotate=90,above]{$l_\alpha=3$}([xshift=-1mm]A-4-1.south west);
  \end{tikzpicture}
  \end{Example}

  The first diagram shows that most of the right-hand diagram is drawn
  by the~\keyword{\Diagram} command, even though most of the commands in
  the example are \TikZ commands. In particular, the shading of the
  boxes in the \emph{$(1,2)$-hook}, and the \emph{$(1,2)$-rim hook}, is
  done by the \option{ribbons} key. (In fact, the right-hand picture
  could be drawn without using a \keyword{tikzpicture} environment by
  using the \option{tikz after} key to add the three
  \keyword{\draw}-commands.)

  \begin{warning}
     Unlike the \TikZ-commands, you should \emph{not} put a semicolon
     after an \aTableau-command when it is used inside a
     \keyword{tikzpicture} environment. If you accidentally add a
     semicolon then, most of the time it will not cause any harm and
     \TikZ will just give you a warning.
  \end{warning}

  The top of the last example defines two \TikZ-styles, \textsf{B} and
  \text{G}, which add the \textsf{b}lue and \textsf{g}reen coloured
  hatchings of the two hooks. Inside the \keyword{tikzpicture}
  environment, the Young diagram is placed at position $(0,0)$ by the
  \keyword{\Diagram} command. The three \keyword{\draw} commands are
  used to add the three labelled arrows to the picture. Please consult
  the very readable \TikZ manual for more details on the \TikZ-commands
  in this example.

  From the viewpoint of \aTableau, the most important point of the last
  example is that the first \keyword{\Diagram} command is a standalone
  diagram, whereas the second one is part of a \keyword{tikzpicture}
  environment because it was given the Cartesian coordinate $(0,0)$. The
  second important point is that this example uses four \emph{named
  coordinates}: (\textsf{A-1-1}), (\textsf{A-1-2}), (\textsf{A-4-1}),
  and (\textsf{A-4-2}). When \aTableau draws a tableau, or a diagram, it
  creates a \TikZ named coordinate (\textsf{A-r-c}) for every box in the
  diagram, where (\textsf{A-r-c}) is the box in row~\textsf{r} and
  column~\textsf{c}. These named coordinates, or anchors, are
  \textit{only} defined for the boxes in the tableau.  Giving finer
  control, the following standard \TikZ \textit{anchors} are available,
  depending on which tableaux convention you are using:

  \begin{equation}\label{E:NodeNames}
    \begin{tikzpicture}[every node/.style={font=\tiny, text=gray}]
        \Diagram(0,0)[scale=4]{1}
        \node[label={above:(A-r-c)}] at (A-1-1){$\color{red}\bullet$};
        \node[label={above:(A-r-c.north)}] at (A-1-1.north){$\color{red}\bullet$};
        \node[label={below:(A-r-c.south)}] at (A-1-1.south){$\color{red}\bullet$};
        \node[label={right:(A-r-c.east)}] at (A-1-1.east){$\color{red}\bullet$};
        \node[label={left:(A-r-c.west)}] at (A-1-1.west){$\color{red}\bullet$};
        \node[label={above right:(A-r-c.north east)}] at (A-1-1.north east){$\color{red}\bullet$};
        \node[label={above left:(A-r-c.north west)}] at (A-1-1.north west){$\color{red}\bullet$};
        \node[label={below right:(A-r-c.south east)}] at (A-1-1.south east){$\color{red}\bullet$};
        \node[label={below left:(A-r-c.south west)}] at (A-1-1.south west){$\color{red}\bullet$};
    \end{tikzpicture}
    \qquad
    \begin{tikzpicture}[every node/.style={font=\tiny, text=gray}]
        \Diagram(0,0)[ukrainian, scale=4]{1}
        \node[label={above:(A-r-c)}] at (A-1-1){$\color{red}\bullet$};
        \node[label={above:(A-r-c.north)}] at (A-1-1.north){$\color{red}\bullet$};
        \node[label={below:(A-r-c.south)}] at (A-1-1.south){$\color{red}\bullet$};
        \node[label={right:(A-r-c.east)}] at (A-1-1.east){$\color{red}\bullet$};
        \node[label={left:(A-r-c.west)}] at (A-1-1.west){$\color{red}\bullet$};
        \node[label={above right:(A-r-c.north east)}] at (A-1-1.north east){$\color{red}\bullet$};
        \node[label={above left:(A-r-c.north west)}] at (A-1-1.north west){$\color{red}\bullet$};
        \node[label={below right:(A-r-c.south east)}] at (A-1-1.south east){$\color{red}\bullet$};
        \node[label={below left:(A-r-c.south west)}] at (A-1-1.south west){$\color{red}\bullet$};
    \end{tikzpicture}
  \end{equation}

  You can change the prefix \textsf{A} in these anchors using the
  \option{name} key. Similar anchors are created for abacus beads.
  These \textit{anchors} are a standard part of \TikZ.  If you change
  the shape of the node, then the available anchors may change; for more
  information see the \TikZ manual.

%  The previous example also shows that the \keyword{\Diagram} command
%  accepts a \option{ribbons} option, which adds one or more
%  \textit{ribbons} to the diagram. Both the ribbon and the individual
%  boxes in the ribbon can be given \TikZ styling; see
%  \autoref{S:RibbonTableaux}. All of the options to these commands are
%  described in the following sections.

  In the example above, the Cartesian coordinate $(0,0)$ is necessary
  because the \keyword{\Diagram} command is used inside a
  \keyword{tikzpicture} environment, but the choice of coordinate is not
  important in the sense that changing the coordinate will not change
  the picture. Cartesian coordinates become more meaningful when the
  \keyword{tikzpicture} environment has more components.  In the next
  example, the key \option{name=B} makes the named
  coordinates for the second diagram take the form (\textsf{B-r-c}),
  instead of the default node names (\textsf{A-r-c}), which are used in
  the first diagram. This allows us to draw (wavy) lines between boxes
  in the two tableaux.

  \begin{example}[lefthand width=0.14\textwidth, sidebyside gap=0mm,label={Ex:ABnames}]
  % \usetikzlibrary{decorations.pathmorphing}
  \begin{tikzpicture}[wave/.style={thick,red,->,decorate,decoration=snake}]
    \aTabset{ribbon style={fill=blue!20, opacity=0.4}}
    \Diagram(0,0)[e=2,entries=residues,ribbons={13crc}]{3,2}
    \Diagram(0,-2.5)[e=2,entries=residues,ribbons={12rcr},name=B]{2^2,1}
    \draw[wave]([shift={(0,-0.05)}]A-2-1.south east)
                --([shift={(0,0.05)}]B-1-1.north east);
  \end{tikzpicture}
  \end{example}

  In this example, the entries in the two tableaux were added by the key
  \option{entries=residues}, which automatically puts the
  \option{e}-residue into each of the boxes in the two tableaux. The
  shading is done by the \option{ribbons} commands.  Notice that
  \option{ribbon style}, which sets the \TikZ-style of the ribbons,
  includes \keyword{opacity=0.4}.  This necessary because the ribbons
  are drawn \textit{after} the tableau entries; see \autoref{SS:order}.
  Alternatively, these tableau entries could be included in the ribbon
  specifications, as described in \autoref{S:RibbonTableaux}.

  The final example in the introduction uses \keyword{\ShiftedTableau}
  inside a \emph{matrix of nodes}:

  \begin{Example}
  % \usetikzlibrary{matrix}
  \begin{tikzpicture}[scale=0.8, arr/.style={->,SteelBlue, thick}]
    \matrix (M)[matrix of nodes,row sep=4mm,column sep=4mm]{
        &    &    & \ShiftedTableau{1234,567} \\
        &    & \ShiftedTableau{1235,467}
        &    & \ShiftedTableau{1243,567}\\
        & \ShiftedTableau{1236,457}
        &    & \ShiftedTableau{1245,367}
        &    & \ShiftedTableau{1253,467}\\
      \ShiftedTableau{1237,456}
        &    & \ShiftedTableau{1246,357}
        &    & \ShiftedTableau{1254,367}\\
        & \ShiftedTableau{1247,356}\\
     };
     \draw[arr](M-1-4)--node[above]{$\sigma_4$}(M-2-3);
     \draw[arr](M-1-4)--node[above]{$\sigma_3$}(M-2-5);
     \draw[arr](M-2-3)--node[above]{$\sigma_5$}(M-3-2);
     \draw[arr](M-2-3)--node[above]{$\sigma_3$}(M-3-4);
     \draw[arr](M-2-5)--node[above]{$\sigma_4$}(M-3-6);
     \draw[arr](M-3-2)--node[above]{$\sigma_6$}(M-4-1);
     \draw[arr](M-3-2)--node[above]{$\sigma_3$}(M-4-3);
     \draw[arr](M-3-4)--node[above]{$\sigma_5$}(M-4-3);
     \draw[arr](M-3-4)--node[above]{$\sigma_4$}(M-4-5);
     \draw[arr](M-3-6)--node[above]{$\sigma_3$}(M-4-5);
     \draw[arr](M-4-1)--node[above]{$\sigma_3$}(M-5-2);
     \draw[arr](M-4-3)--node[above]{$\sigma_6$}(M-5-2);
  \end{tikzpicture}
  \end{Example}

  Notice that even though the \keyword{\SkewTableau} commands are being
  used inside a \keyword{tikzpicture} environment they do not need
  coordinates because the skew tableaux are being placed by the \TikZ
  \keyword{\matrix} command. That is, the \keyword{\SkewTableau}
  commands are not explicit components of the \keyword{tikzpicture}
  environment as hey are all (implicitly) inside \TikZ-nodes.

  The following sections describe the commands defined by the \aTableau
  package and how to use them.  We start with the \keyword{\Tableau}
  command because all of the diagram and tableau commands, except for
  \keyword{\RibbonTableau}, are derived from \keyword{\Tableau}.


  \begin{warning}
    The \aTableau package relies heavily on \ctan{LaTeX3}, so it
    requires an update-to-date version of \LaTeX. You may need to update
    your \LaTeX\ installation if it is older than \TeX Live 2024.
  \end{warning}

  % --------------------------------------------------------------
  \section{Tableaux and Young diagrams}

    \emph{Partitions} are weakly decreasing sequences of non-negative
    integers. They are fundamental objects in algebraic combinatorics
    and representation theory.\index{partition}  Rather than working
    with sequences of integers, it is often useful to visualise
    partitions as Young \emph{diagrams}, which are arrays of
    \emph{boxes}, or \emph{nodes}, in the plane\footnote{In this manual,
    we normally refer to the entries of tableaux and diagrams as
    \textit{boxes}, rather than \textit{nodes}. This is to avoid any
    ambiguity with \TikZ-nodes.}. A \emph{tableau} is a (Young) diagram
    where the boxes are labelled using some alphabet. Equivalently, a
    diagram is a tableau with empty labels.
    \index{young diagram@Young diagram|see{diagram}}
    \index{diagram}
    \index{tableau}
    \index{box}
    \index{node}
    \index{node|see{box}}
    \index{exponential notation}
    \index{partition!exponential notation}

    For example, $\lambda=(4,3,3,2,0,0,\dots)$ is a partition of $12$,
    since the entries sum to~$12$. It is customary to omit the zeros
    when writing partitions and to use exponents for repeated parts, so
    we can write this partition in \emph{exponential notation} as
    $\lambda=(4,3^2,2)$. Many authors identity a
    partition~$\lambda=(\lambda_1,\lambda_2,\dots)$ with its
    \emph{diagram}, which is the set $[\lambda] = \set{(r,c)|1\le
    c\le\lambda_r\text{ for }r\ge1}$.  When using the \emph{English
    convention}, the diagram of~$\lambda$ is visualised as a left
    justified array of boxes in the plane, where the first row has~$4$
    boxes, the next two lower rows have~$3$ boxes and the last row
    has~$2$ boxes. The \emph{shape} of a diagram is the partition that
    gives the number of boxes in each row. Diagrams and tableaux drawn
    using the \emph{French} convention are given by reflecting in a
    horizontal line above the diagrams, whereas the \emph{Ukrainian} and
    \emph{Australian} conventions are given by appropriately rotating
    these diagrams. All of these conventions are used in the literature,
    with some papers using more than one convention.

  % --------------------------------------------------------------
    \subsection{Tableaux}
    \indexcmd{Tableau}
    \label{S:Tableau}
    This section describes the \keyword{\Tableau} command, which draws
    tableaux or labelled diagrams. The syntax of this command is:

    \keyword{\Tableau (x,y) [options] {tableau entries}}

    where:

    \begin{options}
      \item[*$(x,y)$]\index{xy@$(x,y)$} The Cartesian coordinates $(x,y)$
      are needed if, and only if, the tableau is a component of a
      \keyword{tikzpicture} environment.
      \item[*options] Optional arguments. Described, with examples, below.
      \item[*tableau entries]\index{tableau entries} The
      \keyword{tableau entries} are given as a comma-separated list, with commas
      separating rows and each \textit{token}, or braced-group of
      tokens, being in a separate column. As we explain below, each
      tableau entry can be prefixed by (optional) \TikZ-styling
      specifications.
    \end{options}

    We show by example how to use the \keyword{\Tableau} command,
    starting with basic usage and finishing with style. Inside
    \keyword{\Tableau}\!, the rows are separated by commas, with the
    columns given by the ``letters'' in the word for each row:

  \begin{example}
  \Tableau{ 1234, 56, 7, 8 }       \qquad
  \Tableau{ \alpha\gamma\delta q,
            xy\eta,
            \lambda\mu,
            \pi\sum }
  \end{example}

  The tableau specifications can be put on a single line, without any
  spaces. Alternatively, as shown here, spaces and line breaks can be
  added to improve readability. The one restriction is that you cannot
  have blank lines inside a \keyword{\Tableau} command. As this example
  suggests, by default, the entries in a tableau are typeset in
  mathematics-mode, but this can be changed using the
  \option{text entries} key, which is described below.

  As the examples above show, the tableau entries can be individual
  characters, or they can be \TeX\ commands. Tableaux frequently contain
  entries that are not single characters, or given by commands. Such
  entries can be added to a tableau by enclosing them in braces:

  \begin{example}
    \Tableau{ 13{10}{11}{2x}, 2{x^2}, {r_1} }
  \end{example}

  % There is one additional caveat for braced-entries: any
  % \textit{singleton} braced entry \keyword{{...}} in a row needs to be
  % \textit{double-braced}. Double-bracing is only necessary when you have
  % a row that has only one column, and you have a multi-letter column
  % entry. The need for double-bracing is a consequence of how \LaTeX3
  % detects braced commas in comma-separated lists.
  % \index{tableau entries!double-braced}
  % \index{double-braced}

  % The following example shows the difference between single and
  % double-bracing for tableau entries.
  % \index{tableau entries!double-braced}

  % \begin{example}
  % \Tableau{ 1345, 8{10}, {11}, {12} } \qquad
  % \Tableau{ 1345, 8{10}, {{11}}, {{12}} }
  % \end{example}

  % The $10$'s in these tableaux do not need to be doubled-braced because
  % they appear in a row of length~$2$, whichi s greater than~$1$!. In
  % contrast,  because they are not doubled-braced, the~$11$ and~$12$ in
  % the first tableau are both treated as being two characters in
  % different columns in the first tableau. In the second tableau,
  % the~$11$ and~$12$ each appear in a separate column because they are
  % double-braced.

  % --------------------------------------------------------------
  \subsubsection{Adding style}
  \index{star}
  \index{*|see{star}}
  \index{tableau entries!style}
  \index{tableau entries!star}

  The \aTableau package provides several different ways to add \TikZ
  style specifications to the boxes in diagrams and tableaux.  This is
  done by adding style prefixes to the tableau entries.  To take
  advantage of these styles you need to have at least rudimentary
  knowledge of \TikZ-styling. The examples in this manual might be
  enough to get you started as they give a good indication of what is
  possible.  For anything more exotic, please consult the \TikZ manual.

  The simplest way to change the style of an entry in a tableau is to
  add a $*$-prefix to the entry:

  \begin{example}
  \Tableau{ 1*2*3*4*5, 6*789, {10}*{11}{12}, {13} }
  \end{example}

  In this tableau, there is an asterisk before each of the entries $2$,
  $3$, $4$, $5$ and $8$, which gives the corresponding boxes the
  \option{star style}. The $*$-entries are given the
  \TikZ-styling, which is  \keyword{fill=}\aTableauColour{StarStyle} by default. You can
  change the default $*$-style using the \option{star style} key:

  \begin{example}
  \Tableau[star style={fill=red!20,draw=red,thick}]
          { 1*2*3*4*5, 6*789, {10}*{11}{12}, {13} }
  \end{example}

  Note that the $*$-styling overrides the default styling of the
  surrounding boxes. This is because the boxes with the default styling
  are placed first, in the order that they are entered, after which the
  boxes with custom styles are placed.

  The $*$-syntax is a quick shorthand for adding emphasis to some boxes
  in a tableau, but all of the $*$-entries are given the same style.
  It is possible to give every box in the tableau different styles by
  putting \TikZ-styling specifications inside square brackets,
  \keyword{[...]}, before the corresponding letter in the tableau
  specification.  You can mix the $*$-syntax and the
  \keyword{[...]}-style syntax for different entries, although only one
  of these style settings can be applied to any given box.

  \begin{example}
  \Tableau[french]{
      1 [blue!20]2 [circle]3 [red]4 [cyan]5,
      6 *7 8 9,
      {10} *{11} {12},
      {13} }
  \end{example}

  It is not necessary to add spaces between the entries in the way that
  this example does. This is done only to make it clearer how the
  style specifications are applied to the different entries in the tableau.

  Omitting some technicalities, each box in a tableau is constructed as
  a \TikZ node, so any \TikZ-styling specifications for a node can be
  used.  This example shows that some care is needed when changing the
  style because simply giving a colour changes both the colour used to
  \textit{fill} the box and the \textit{text colour}, which is why some
  of tableau entries in the last example appear to be blank.  If you are
  already familiar with \TikZ then you probably already know what to do.
  If not, then here is a short list of useful style specifications for
  \TikZ nodes:
  \begin{center}\index{fill}\index{text}\index{draw}\index{font}
     \begin{tabular}{ll}\toprule
       Style & Meaning \\\midrule
       \keyword{draw}$=\langle\text{colour}\rangle$
            & Sets the boundary colour of the node\\
       \keyword{fill}$=\langle\text{colour}\rangle$
            & Sets the fill, or background, colour of the node\\
       \keyword{font}$=\langle\text{font commands}\rangle$
            & Sets the font \\
       \keyword{opacity}$=\langle\text{value}\rangle$
            & Sets both the drawing and filling opacity to $=\langle\text{value}\rangle$\\
       \keyword{text}$=\langle\text{colour}\rangle$
            & Sets the text colour in the node \\
       \bottomrule
     \end{tabular}
  \end{center}
  In the style settings, any valid \LaTeX\ colour name can be used; see, for
  example, the \ctan{xcolor} manual. In addition, the style specifications
  \keyword{thin}\!,
  \keyword{thick}\!,
  \keyword{very thick}\!,
  \keyword{ultra thick}\!,
  \keyword{dashed}\!,
  \keyword{dotted}\!, $\dots$.
  change the border thickness, and its properties.
  \autoref{SS:TableauOptions} below describes how to use these options
  to customise the tableau style.

  If you want to apply ``complicated'' style settings, such as
  \keyword{draw=cyan,ultra thick}, which consists of a comma-separated
  list of $\TikZ-$style settings, then the entire style must be put
  inside matching square brackets and curly braces \keyword{[{...}]}.
  This is necessary to ensure that  \LaTeX\ does not get confused by the
  commas in the style setting (since commas are also used to separate
  the rows of the tableau).

  \begin{example}
  \Tableau[australian]{
      1 [fill=blue!20]2 3 [text=red]4 5,
      [{circle,fill=orange}]6 *7 8 9,
      [{draw=cyan,ultra thick}]{10} *{11} {12},
      [{dashed,draw=red,ultra thick}]{13} }
  \end{example}

  % The last example shows that it is not necessary to double-brace
  % singleton entries when they have a style specification.

  This example shows that modifying the borders of a box is problematic
  because boxes placed later are likely to overwrite the new border
  style (see \autoref{SS:order}). As a result, modifying the borders of
  the boxes is often not a good way to add emphasis to a tableau,  so
  use this sparingly. In addition, the orange circle in this example
  looks too big, but it is what we asked for because the diameter of the
  circle is the horizontal width of the box. The next example shows that
  we get a better result with \textit{minimum size} specification when
  using the \textit{Australian} and \textit{Ukrainian} conventions  for
  tableaux.

  For complex or frequently used styles, we recommend using the
  \keyword{\tikzset} command to define the style outside of the
  \keyword{\Tableau} command. This makes your code both easier to read
  and easier to change. If you need to define a style for a single
  picture, then you can also use the \option{styles} key.

  \begin{example}[lefthand width=0.3\textwidth]
  \tikzset{
     B/.style={fill=blue!40, text=yellow},
     G/.style={fill=green!80!blue,circle,minimum size=5mm},
     R/.style={fill=red,text=white, font=\tiny},
     Y/.style={fill=yellow!20, text=blue}
  }
  \Tableau[french]{[R]1[B]2,[G]3,[Y]4} \quad
  \Tableau[ukrainian]{[R]1[B]2,[G]3,[Y]4} \newline
  \Tableau[english]{[R]1[B]2,[G]3,[Y]4} \quad
  \Tableau[australian]{[R]1[B]2,[G]3,[Y]4}
  \end{example}

  In this example, the \TikZ-styles \textsf{B}, \textsf{G} \textsf{R},
  and \textsf{Y} are applied to the tableaux entries labelled $1$, $2$,
  $3$, and $4$, respectively. Note the use of \emph{minimum size} to
  control the size of the circle. The colours used here, and throughout
  this manual, are for demonstration purposes only. I recommend against
  such garish choices in any self-respecting document!

  % --------------------------------------------------------------
  \subsubsection{Tableau coordinates}\label{SS:coordinates}
  \index{xy@$(x,y)$}
  \index{Cartesian coordinates!tableau}

  The $(x,y)$-coordinates are required if, and only if, the
  \keyword{\Tableau}command is used inside a \keyword{tikzpicture}
  environment, in which case $(x,y)$ gives the coordinates of the
  ``outside corner'' of the box in row~$1$ and column~$1$ of the
  tableau. The following example shows how tableaux are placed using
  $(x,y)$-coordinates inside a \keyword{tikzpicture} environment. The example
  also shows that, by default, the tableau boxes are square and half a unit wide.

  \tikzset{
    add grid/.code = {
      % draw coordinate axes
      \draw[help lines, step=0.5,thin, LightGrey](-0.5,-1.5) grid (4,2.5);
      \draw[help lines, step=1,thin, Grey](-0.5,-1.5) grid (4,2.5);
      \foreach \pt in {0,...,3} {
        \node[below,font=\tiny] at (\pt,-1.5){$\pt$};
      }
      \foreach \pt in {-1,...,2} {
        \node[left, font=\tiny] at (-0.5,\pt){$\pt$};
      }
    }
  }

  \begin{example}
  \begin{tikzpicture}[add grid]
    \Tableau(0,0)    [english]   {123,45}
    \Tableau(0,0.5)  [french]    {123,45}
    \Tableau(2.5,0.5)[ukrainian] {123,45}
    \Tableau(2.5,0)  [australian]{123,45}
  \end{tikzpicture}
  \end{example}

  % --------------------------------------------------------------
  \subsubsection{Tableau keys}
  \label{SS:TableauOptions}

  The optional $*$-styles and $[\ldots]$-styles are the main mechanism
  for styling the individual boxes in a tableau. Using a key-value
  syntax, you can change aspects of the tableau produced by the
  \keyword{\Tableau} command.  The rest of this section describes these
  keys, their default values, and gives examples of how to use them.

  The options below can be applied to the \keyword{\Tableau}
  command by giving them as a comma-separated list inside square
  brackets:
  \begin{center}
    \keyword{\Tableau [options] {tableau entries}}\quad or \quad
    \keyword{\Tableau (x,y) [options] {tableau entries}}.
  \end{center}
  The options can appear in any order, with later options having
  precedence over earlier ones.

  When used inside a \keyword{\Tableau} command, the options only affect
  that particular tableau. Most of these options, or keys, can also be
  used in the \keyword{\aTabset} command, to change the default settings
  of all tableaux in the same \LaTeX\ group, or as optional arguments in
  \keyword{\usepackage[..]{atableau}}, to change the default settings
  for the entire document.  For example, if you put the command
  \keyword{\aTabset[ukrainian]} in the preamble of your document then,
  by default, all subsequent tableaux and Young diagrams will be drawn
  using the Ukrainian convention.

  \option[-]{english}
  \option*[]{french}
  \option*[]{ukrainian}
  \option*[]{australian}

  These four (mutually exclusive) options change the convention, or
  orientation, that is used when drawing the tableaux. Rather than
  defining these conventions precisely, we use the tired and true method
  in combinatorics of defining by example.

  \begin{example}[lefthand width=0.3\textwidth]
  \Tableau[ukrainian, styles={R={fill=red!20,draw=red}}]
    {12[R]3,45}
  \end{example}

  \begin{example}[lefthand width=0.3\textwidth]
  \tikzset{R/.style={fill=red!20, circle, draw=red, ultra thick}}
  \Tableau[french]{12[R]3,45}
  \end{example}

  \begin{example}[lefthand width=0.3\textwidth]
  \tikzset{B/.style={fill=blue!20, dashed, thick}}
  \Tableau[english]{12[B]3,45}
  \end{example}

  \begin{example}[lefthand width=0.3\textwidth]
  \Tableau[australian, styles={B={fill=blue!20, text=white}}]
    {12[B]3,45}
  \end{example}

  Changing the convention sets the \option{box height} and \option{box
  width}, relative to the current \option{scale}. If you want to use a
  custom box height and box width, then you need to set them
  \emph{after} specifying the convention.

  By default, the \option{english} convention is used to draw tableaux
  and Young diagrams. The default convention can be changed using
  \keyword{\aTabset}. Capitalised names, \keyword{Australian}\!,
  \keyword{English}\!, \keyword{French} and \keyword{Ukrainian}\!, are also
  recognised. If, for example, your document is only going to draw
  \option{french} tableaux and diagrams, then you can specify this when
  you load the package using \keyword{\usepackage[french]{atableau}}\!,
  or by adding \keyword{\aTabset{french}} to your document preamble.

  \option[centre]{align=value}[centre/bottom/top]

  By default, the baseline of a tableau is its centre. This can be
  changed using the \option{align} key.

  \begin{example}
    % the default
    Some \Tableau[align=centre]{124,357,68} words
  \end{example}

  \begin{example}
    Some \Tableau[align=top]{124,357,68} words
  \end{example}

  \begin{example}
   Some \Tableau[align=bottom]{124,357,68} words
  \end{example}

  Similarly, use \option{align} to control how tableaux are aligned
  in displayed equations.

  \begin{example}
    \aTabset{align=bottom}
    \Tableau{124,357,68}\qquad
    \Tableau{124,35}
  \end{example}

  \begin{example}
    % by default: align=centre
    \Tableau{124,357,68}\qquad
    \Tableau{124,35}
  \end{example}

  \begin{example}
    \aTabset{align=top}
    \Tableau{124,357,68}\qquad
    \Tableau{124,35}
  \end{example}

  \option[true]{border}[true/false] \newline
  \option[false]{no border}[false/true]

  The \option{border} and \option{no border} keys enable and disable The
  \option{no border} key is inverse to \option{border}, so \option{no
  border} is equivalent to \option{no border=true}the drawing of the
  (thick) border wall around a tableau.

  \begin{example}
  \Tableau[ukrainian, border=false]{123,455,674} \qquad
  \Tableau[french, no border]{123,455,674} \qquad
  \end{example}

  \option[\aTableauColour{Main}]{border colour=colour}[a \LaTeX~colour]
  \option*[thick, line~cap=rect]{border style=style}[\TikZ-styling]

  The \option{border colour} option sets the colour of the border wall,
  whereas \option{border style} sets the \TikZ-styling of the border wall.
  (So, \option{border style} can override \option{border colour}.)

  \begin{example}
  \Tableau[ukrainian, border colour=red]{123,455,674} \qquad
  \Tableau[french, border style={dashed,orange}]{123,455}
  \end{example}

  The \option{border style} key \emph{appends} any \TikZ-styles to the
  current styling of the wall.

  % nodes
  \option[none]{box fill=colour}[a \LaTeX~colour]

  The \option{box fill} option sets the background colour of a box in a
  tableau, which is used in exactly the same way as the \textsf{fill} key for a \TikZ-node. If you use
  a dark fill colour, then you will almost certainly want to change the
  text colour using the \option{box text} key -- and you may also want to
  change the colours of the border walls and inner walls using
  \option{border colour} and \option{inner wall}, respectively.

  \begin{example}
  \Tableau[box fill=blue, box text=yellow]{123,455,674} \qquad
  \Tableau[box fill=red!10]{123,455}
  \end{example}

  \option[]{box font=font command}[\LaTeX\ font command]

  The \option{box font} key sets the font used to type the entries of a
  tableau.  By default, tableau entries are typeset as mathematics, so
  this is mainly useful for changing the font size (because, for example
  \keyword{\bfseries $1$} does not make the $1$ bold). It is only when
  you are using \option{text entries} that font commands like
  \keyword{\itshape} and \keyword{\bfseries} will have any effect.

  \begin{example}
  \Tableau[box font=\tiny]{123,455,674} \qquad
  \Tableau[text entries, box font=\bfseries]{123,455}
  \end{example}

  \option[0.5]{box height=number}[a decimal giving the height in \unit{cm}]
  \option*[0.5]{box width=number}[a decimal giving the width in \unit{cm}]

  Sets the width and height of the boxes in the tableau. The default
  height and width of the boxes is \qty{0.5}{cm} when using the
  \option{english} and \option{french} conventions and
  \qty{0.7012}{cm}, which is approximately $1/\sqrt{2}$\unit{cm}, when
  using the \option{ukrainian} and \option{australian} conventions. (In
  all cases, the default side length of the boxes is \qty{0.5}{cm}.)

  \begin{example}
  \aTabset{align=top}
  \Tableau[box height=0.8]{123,455,674} \qquad
  \Tableau[box width=0.2, ukrainian]{123,455}
  \end{example}

  The \option{box height} and \option{box width} keys need to be set
  \emph{after} changing the tableau convention using
  \option{australian}, \option{english}, \option{french}, or
  \option{ukrainian} because the conventions change the height and width
  of the tableau boxes (and later conventions override earlier ones). In
  particular, in the last example, \option{box width=0.2} has no effect
  because following \option{ukrainian} key sets the box width to
  \qty{0.7012}{cm}.

  \begin{example}
  \aTabset{align=top}
  \Tableau[box height=0.3]{123,455,674} \qquad
  \Tableau[ukrainian, box width=0.3]{123,455}
  \end{example}

  See also the \option{scale}, \option{xscale} and \option{yscale}
  options.

  \option[\aTableauColour{Main}]{box text=colour}[a \LaTeX~colour]

  As for a \TikZ-nodes, use \option{box text} to set the default text
  colour for the entries of the tableau boxes.

  \begin{example}
  \Tableau[box text=red]{123,45} \qquad
  \Tableau[box text=blue, french]{123,45}
  \end{example}

  \option[]{box style=style}[\TikZ~styling]

  All of the preceding keys for tableau boxes can be overridden by using
  \option{box style} to set the style of the tableau boxes directly.

  \begin{example}
  \Tableau[box style={text=red, draw=orange,shape=circle}]{123,45} \qquad
  \Tableau[box style={font=\small}, french]{123,45}
  \end{example}

  \option[false]{conjugate}[false/true]

  Draws the \emph{conjugate} tableau, where the rows and columns are
  swapped.

  \begin{example}
  \Tableau{123,45,67}    \qquad
  \Tableau[conjugate]{123,45,67}
  \end{example}

  \option[centre]{halign=value}[centre, left, right]
  \option*[centre]{valign=value}[bottom, centre, top]

  By default, the entries in tableaux are centered, both horizontally
  and vertically, and it is your responsibility to ensure that the
  entries fit inside the tableau boxes (you can change the box
  dimensions using the \option{box height}, \option{box width} and
  \option{scale} keys). The \option{halign} and \option{valign} keys
  provide a crude way to change the horizontal and vertical alignment of
  the box entries.

  \begin{example}
  \Tableau[halign=left]{1{11}{11},{11}1} \quad
  \Tableau[halign=centre]{1{11}{11},{11}1} \quad
  \Tableau[halign=right]{1{11}{11},{11}1}
  \end{example}

  \begin{example}
  \Tableau[valign=bottom]{qf,gpt,jb} \quad
  \Tableau[valign=centre]{qf,gpt,jb} \quad
  \Tableau[valign=top]{qf,gpt,jb}
  \end{example}

  These keys are provided for completeness only. We recommend not using them!

  \option[\aTableauColour{Inner}]{inner wall=colour}[\LaTeX~colour]
  \option*[thin, line cap=rect]{inner style=style}[\TikZ-styling]

  The \option{inner wall} option sets the colour of the inner wall,
  whereas \option{inner style} sets the \TikZ-styling of the inner wall.
  (In particular, \option{inner style} can be used to override
  \option{inner wall}.)

  \begin{example}[label=Ex:dashed]
  \Tableau[australian, inner wall=red]{123,455,674} \qquad
  \Tableau[inner style=dashed]{123,455}
  \end{example}


  \index{dashed}
  \index{style!dashed}
  If you look closely at the second tableau in \autoref{Ex:dashed}, you will
  notice that the dashes are not very clean. This is because, for
  example, the dashes on bottom of one box do not properly match up with
  the dashes on the top of the adjacent box below it. Rather than
  using \option{inner style=dashes} it is better to use something like
  \option{inner style=\{dash pattern=on 1pt off 2pt\}}.

  \begin{example}
  \Tableau[inner style={dash pattern=on 1pt off 2pt}]
    {123,455}
  \end{example}

  You can use \option{inner wall=none}, to remove the inner tableau walls. Alternatively, use \option{no boxes}.

  \begin{example}
  \Tableau[ukrainian, inner wall=none]{123,45,67} \qquad
  \Tableau[ukrainian, no boxes]{123,45,67}
  \end{example}

  \option[none]{label=text}[any character/text]
  \option*[font=\textbackslash scriptsize, text=\aTableauColour{Inner}]{label style=style}[\TikZ-styling]

  The \option{label} key adds a label to a tableau next to the
  $(1,1)$-box. The style of the label can
  be changed using the \option{label style} key.

  \begin{example}
  \Tableau[ukrainian, label={\mathtt{t}}]{13,25,4} \quad
  \Tableau[french, label=0, label style={red}]{13,5} \quad
  \Tableau[label style={draw=orange,circle, inner sep=0pt,
    minimum size=2mm}, label=1] {12,35,5}
  \end{example}

  Like tableau entries, \option{label} is typeset as mathematics by default .

  \option[true]{math entries}[true/false]
  \option*[false]{text entries}[true/false]

  Use the \option{math entries} and \option{text entries} keys to have
  the tableau entries typeset as either mathematics or text,
  respectively. By default, the entries of a tableau are typeset as
  mathematics.

  \begin{example}
  \Tableau[math entries]{ABCD, efg, HI}\qquad% the default
  \Tableau[text entries]{ABCD, efg, HI}
  \end{example}

  \option[A]{name=text}[string]

  By default, the boxes can be referenced using the node names
  \textsf{(A-1-1)}, \textsf{(A-1-2)}, \textsf{(A-1-3)}, \dots, with
  \textsf{(A-$r$-$c$)} referring to the box in row~$r$ and column~$c$.
  Use \option{name} to change the \emph{prefix} \textsf{A} to
  anything else that you like.  This feature is most useful when you
  have several tableaux inside a \keyword{tikzpicture} environment
  because by using \option{name} you can refer to boxes in the different
  diagrams. If you only have one tableau, then you would not normally
  need to change the default prefix for the node names. See
  \autoref{E:NodeNames} for a description of the extended anchor names.
  For examples using node names and anchors see \autoref{Ex:ABnames},
  \autoref{Ex:pointing} and \autoref{Ex:waves}.

  \option[]{paths=path specifications}[list of paths]
  \option*[]{path style=style}[\TikZ-styling]

  Use the \option{paths} key to add a comma-separated list of
  \emph{ribbon paths} to the tableau. This allows you to draw certain
  types of paths inside your tableaux and diagrams.

  Ribbon paths are specified using the \keyword{\RibbonTableau} syntax,
  where you first specify the row and column indices of the \emph{head}
  of the ribbon, which is the unique node of maximal content in the
  ribbon, and then give a sequence of \textsf{r}'s and \textsf{c}'s
  depending on whether the row index increases, or the column index
  decreases, respectively.  As always, you have the option of adding
  style --- and you can also specify the contents of each box in the
  ribbon.  For more details about the ribbon specifications, with
  examples, see \autoref{S:RibbonTableaux}, which describes the
  \keyword{\RibbonTableau} command.

  \begin{example}
  \Tableau[styles={R={red,thick}, C={cyan,thick}},
     paths={(R)14ccrc,(C)23rcc}]{1234,567,89{10}}
  \end{example}

  The \option{path style} key changes the default style of the path:

  \begin{example}
  \Tableau[styles={R={red,thick}, C={cyan,thick}},
     paths={(R)14ccrc,(C)23rcc},
     path style={rounded corners}]{1234,567,89{10}}
  \end{example}

  \option[]{path box=text}[a node entry]
  \option*[]{path box style=style}[\TikZ-styling]

  The \option{path box} sets the default entry for every box on a path.
  The \option{path box} is only used if you have not specified a box
  entry as a subscript in the path/ribbon specification.  This key is
  useful if you want to mark all of the boxes in every path with the
  same symbol.

  \begin{example}
  \tikzset{R/.style={red,thick},C/.style={cyan,thick}}
  \Tableau[paths={(R)14ccrc,(C)23rcc},
    path box=\bullet
  ]{1234,567,89{10}}
  \end{example}

  Use \option{path box style} to change the style of the boxes in
  the path.

  \begin{example}
  \Tableau[styles={R={red,thick}, C={cyan,thick}},
    paths={(R)14ccrc,(C)23rcc},
    path box style={draw,circle, minimum size=4mm}
  ]{1234,567,89{10}}
  \end{example}

  By combining \option{path box style} and \option{path style}, you can
  ``decorate'' the boxes on a path.

  \begin{example}
  \Tableau[styles={R={red,thick}, C={cyan,thick}},
    paths={(R)14ccrc,(C)23rcc},
    path box style={draw,circle, minimum size=4mm},
    path style={draw=none},
  ]{1234,567,89{10}}
  \end{example}

  Another way to do this is to change the style of the individual paths,
  which gives more control because different paths can be given
  different styling.

  \begin{example}
  \Tableau[styles={R={draw=none,text=red,thick},
                   C={cyan,thick}},
    paths={(R)14ccrc,(C)23rcc},
    path box style={draw,circle, minimum size=4mm},
  ]{1234,567,89{10}}
  \end{example}

  \option[]{ribbons=ribbon specifications}[list of ribbons]
  \option*[draw=\aTableauColour{Inner}, thin]{ribbon style=style}[\TikZ-styling]
  \option*[]{ribbon box=text}[A node entry]
  \option*[]{ribbon box style=style}[\TikZ-styling]

  Use the \option{ribbons} key to add a comma-separated list of
  \emph{ribbons} to a tableau.  The ribbon specifications are identical
  to those used for \option{paths}.  For more details see
  \autoref{S:RibbonTableaux}, which describes the
  \keyword{\RibbonTableau} command.

  The \option{ribbons} key works in almost exactly the same way
  as the \option{paths} key, with the difference being that we are
  adding ribbons to the tableau rather than lines. The
  \option{ribbon style} key controls the default ribbon style. The
  \option{ribbon box} and \option{ribbon box style} keys work in the
  same way as \option{path box} and \option{path box style},
  respectively.

  \TikZ-styling can be added to ribbons using \option{ribbon style} and
  \option{ribbon box style}. The \option{ribbon style} changes the
  default style of the entire ribbon, whereas \option{ribbon box style}
  changes the default style of the individual boxes in the ribbon.

  \begin{example}
  \Tableau[ribbon style={fill=blue!20,opacity=0.4},
    ribbons={22rc}]{123,45,67}
  \end{example}

  Use \option{ribbon box style} to change the default style of the
  boxes in a ribbon. For example, noting that a box is a ribbon of
  length~$1$, the following code uses ribbons of length one (that is,
  nodes), to highlight the addable nodes of this tableau.

  \begin{example}
  \Tableau[star style={fill=green!80!blue,text=white},
    ribbon style={fill=yellow,text=red},
    ribbons={15_A,24_B,42_C,51_D},
  ]{123*4,567,89*{10},*{11}}
  \end{example}

  As this example indicates, ribbons are assumed to be contained inside
  the diagram of the tableau, so they do not change the border of a
  tableau.

  \option[]{snobs=snob specifications}[list snobs]
  \option*[draw=\aTableauColour{Inner}, thin]{snob style=style}[\TikZ-styling]
  \option*[]{snob box=text}[A node entry]
  \option*[]{snob box style=style}[\TikZ-styling]

  Use \option{snobs} to add ribbons to the tableau, which are ribbons
  that are placed \textit{after} the border of the diagram is drawn.
  (Writing \textit{ribbons} backwards gives \textit{snob}bir.) The
  \option{snob style} key controls the default style of snobs.

  \begin{example}
  \Tableau[star style={fill=green!80!blue, text=white},
    snob style={fill=yellow, text=red},
    snobs={15_A,24_B,42_C,51_D},
  ]{123*4,567,89*{10},*{11}}
  \end{example}

  The difference between this pictures and the one given above using
  \option{ribbons} is that in the \option{snobs} are drawn over the top
  of the tableau border, causing the border to disappear for these
  nodes.  For this reason, in most cases you should normally avoid
  \option{snobs} and use \option{ribbons}. This said, \option{snobs} are
  exactly what are needed in \autoref{Ex:Garnir}.

  The \option{snob box} and \option{snob box style} keys work in the
  same way as the \option{path box} and \option{ribbon box}, and
  \option{path box style} and \option{ribbon box style}, keys.

  \option[1]{scale=number}[a decimal number]
  \option*[1]{xscale=number}[a decimal number]
  \option*[1]{yscale=number}[a decimal number]

  By design, boxes in tableaux do not automatically resize to fit
  their contents. For example, consider:

  \begin{example}
  \Tableau{ 13{10}{11}{1+x}, 2{1+x} }
  \end{example}

  This is unlikely to be the desired output! The options \option{scale},
  \option{xscale} and \option{yscale} can be used to rescale the boxes
  in tableaux and diagrams. The names are slightly misleading because
  \option{xscale} rescales in the direction of increasing column index
  and \option{yscale} rescales in the direction of increasing row index.

  \begin{example}
  \Tableau[scale=2]{ 13{10}{11}{1+x}, 2{1+x} }
  \end{example}

  \begin{example}
  \Tableau[xscale=2]{ 13{10}{11}{1+x}, 2{1+x} }
  \end{example}

  When using the \textsf{australian} and \textsf{ukrainian} conventions,
  \option{xscale} and \option{yscale} scale in both the \textsf{x}
  direction and the \textsf{y}-direction, reflecting how these tableaux
  grow with increasing column and row index, respectively.

  \begin{example}
  \Tableau[ukrainian, xscale=2]{ 13{10}{11}{1+x}, 2{1+x} }
  \end{example}

  \begin{example}
  \Tableau[australian, yscale=2]{ 13{10}{11}{1+x},2{1+x} }
  \end{example}

  \begin{warning}
    The options, \option{scale}, \option{xscale}, and \option{yscale},
    affect all \aTableau diagrams, including tableaux, Young diagrams
    and abacuses. If you want to use \keyword{\aTabset} to change the
    default sizes of the diagrams and tableaux in your document, then it
    is probably better to use the \option{box height} and \option{box
    width} options.
  \end{warning}

  \option[0.5]{script=value}[decimal number]
  \option*[0.4]{scriptscript=value}[decimal number]

  Although rare, tableaux and diagrams are sometimes used as subscripts
  and superscripts. The pacakage automatically rescales tableau when
  they are used as subscripts and superscripts, with the \option{script}
  and \option{scriptscript} keys giving finer control over the scaling
  of subscripts and subsubscripts, respectively. The current values of
  the \option{scale}, \option{xscale} and \option{yscale} keys are taken
  into account.

  \begin{example}
    $v_{\Tableau{123,45}}$ \qquad% the default
    $v_{\Tableau[script=0.3, xscale=2]{123,45}}$
  \end{example}

  \option[false]{shifted}[true/false]

  Set to \textsf{true} for a shifted tableau or shifted diagram. Shifted
  tableaux, and shifted diagrams, are discussed in more detail in
  \autoref{S:Shifted}, which describes the \keyword{\ShiftedDiagram} and
  \keyword{\ShiftedTableau} commands.

  \begin{example}
  \Tableau[shifted]{123,45} \qquad
  \Diagram[shifted]{3,2^2}
  \end{example}

  See \autoref{S:Shifted} for the options specific to shifted tableaux and
  shifted diagrams.

  \option[0]{skew=partition}[a partition]

  Specify the skew shape for a skew tableau or skew diagram. Skew
  tableaux, and skew diagrams, are discussed in more detail in
  \autoref{S:Skew}, which describes the \keyword{\SkewDiagram} and
  \keyword{\SkewTableau} commands.

  \begin{example}
  \Tableau[skew={2}]{123,45} \qquad
  \Diagram[skew={2^2,1}]{3,2^2}
  \end{example}

  See \autoref{S:Skew} for the options specific to skew tableaux and skew diagrams.

  \option[fill=\aTableauColour{StarStyle}]{star style=style}[\TikZ-styling]

  Use \option{star style} to change the style of the starred entries
  of a tableau. (In \TikZ parlance, \option{star style} \emph{appends}
  appends these styles to the default \option{star style}.)

  \begin{example}
  \Tableau[star style={text=magenta, draw=cyan, thick}]
    {1*2*3,4*5} \qquad
  \Tableau[star style={fill=orange!50}]{1*2*3,4*5}
  \end{example}

  Note that the tableau border is drawn \textit{after} the
  \option{star style} is applied. You can use \option{snobs} to draw
  over the top of the border.

  \option[]{styles=\TikZ styles}[List of styles]

  The \option{styles} key is a shorthand for defining \TikZ-styles that
  are only used in the current picture. For example, instead of writing

  \begin{example}
  \tikzset{
     B/.style={fill=blue!40, text=yellow},
     G/.style={fill=green!80!blue,circle,minimum size=5mm},
     R/.style={fill=red,text=white, font=\tiny},
     Y/.style={fill=yellow!20, text=blue}
  }
  \Tableau[french]{[R]1[B]2,[G]3,[Y]4}
  \end{example}

  you can save some typing using the \option{styles} key:

  \begin{example}
  \Tableau[french, styles={
     B={fill=blue!40, text=yellow},
     G={fill=green!80!blue,circle,minimum size=5mm},
     R={fill=red,text=white, font=\tiny},
     Y={fill=yellow!20, text=blue},
  }]{[R]1[B]2,[G]3,[Y]4}
  \end{example}

  Styles defined this way will only be available in the current picture.
  Use \keyword{\tikzset} whenever you are defining styles that are likely
  to be used more than once.

  \option[false]{tabloid}[true/false]

  Use the \option{tabloid} key to draw a tabloid. This key is discussed in
  more detail in \autoref{S:Tabloid}, which describes the
  \keyword{\Tabloid} command.

  \begin{example}
  \Tableau[tabloid]{123,45} \qquad
  \Diagram[tabloid]{3,2^2}
  \end{example}

  See \autoref{S:Tabloid} for the options specific to tabloids.

  \option[]{tikz before=commands}[\TikZ commands]
  \option*[]{tikz after=commands}[\TikZ commands]

  These keys inject \TikZ code into the underlying \keyword{tikzpicture}
  environment, where the tableau is constructed, before and after the
  \keyword{\Tableau} command, respectively.

  \begin{example}[lefthand width=0.24\textwidth]
  \Tableau[
   tikz before={\draw[help lines, step=0.5](-0.5,-2)grid(2,0.5);},
   tikz after={\draw[ultra thick, red]
                 (A-1-3.north east)--(A-1-3.south east)
                 --(A-2-1.north east)--(A-3-1.south east)
                 --(A-3-1.south west);
   }]{134,567,89}
  \end{example}

  Both before and after hooks are provided because the order in which
  you place \TikZ commands affects the final picture because later commands
  are drawn over the top of earlier ones. Note that the \emph{before}
  code cannot use the named nodes \textsf{A-r-c} because this code is
  executed before these nodes are constructed, which happens when the
  tableau is drawn. For more examples, see \autoref{SS:order}.

  \option[]{tikzpicture=keys}[\TikZ-keys]

  This command adds an optional argument to the underlying
  \keyword{tikzpicture} environment, which contains the tableau. The
  \option{tikzpicture} option is ignored  when the \keyword{\Tableau}
  command is equipped with $(x,y)$-coordinates.

  As a first example, we can use \keyword{tikzpicture={rotate=30}}
  to rotate a tableau by 90 degrees:

  \begin{example}
  \Tableau[tikzpicture={rotate=90}]{123,45} \qquad
  \begin{tikzpicture}[rotate=90]
   \Tableau(0,0){123,45}
  \end{tikzpicture}
  \end{example}

  The two tableaux constructions in this example are equivalent. An
  oddity with both of these examples is that the numbers in the tableaux
  have not been rotated, which is because the \TikZ \emph{rotate} key
  does not affect component objects, such as nodes, inside a \TikZ
  picture. You can rotate the entire picture using the
  \TikZ-key \textit{transform canvas}

  \vspace*{15mm}

  \begin{example}[valign=top, lefthand width=0.3\textwidth]
  \Tableau[tikzpicture={transform canvas={rotate=90}}]{123,45}
  \end{example}

  A more interesting application is using
  \option{tikzpicture=remember picture} to annotate a tableau:

  \begin{example}[lefthand width=0.3\textwidth, label={Ex:pointing}]
  \Tableau[ukrainian, name=A,
    tikzpicture={remember picture}]
    { 1[fill=blue!20]23[text=red]45,
      6*78[{draw=cyan,ultra thick}]9,
      [{circle,fill=orange}]{10}*{11}{12},
      [{dashed,draw=red,ultra thick}]{13}
    }
  \end{example}

  Since \option{name}=A, the nodes in this tableau are referenced as
  \textsf{(A-1-1)}, \textsf{(A-1-2)}, \textsf{(A-1-3)}, \dots. As the
  example uses \textit{remember picture}, other \keyword{tikzpicture}
  environments can reference the node names in last tableau, which
  allows us to refer back to the table above using:

  \vspace*{3mm}

  \begin{example}[lefthand width=0.3\textwidth]
  \tikz[remember picture, overlay]
   \draw[very thick, ->, red]
     (0,0) node[right,align=left]
        {This circle is too big!\\Use minimum size=5mm}
     to [out=135, in=225](A-3-1);
  \end{example}

  (See the \ctan{tikzmark} package for extensions of this idea.)

  % --------------------------------------------------------------
  \subsubsection{Compositions}
  \index{composition}

  As we have defined them, tableaux are always of \textit{partition}
  shape, in the sense that the lengths of the rows is a weakly
  decreasing sequence. In fact, the \keyword{\Tableau} command also
  accepts tableaux of \emph{composition} shape, where the lengths of the
  rows are not necessarily in decreasing order:

  \begin{example}
  \Tableau{ 45,2751,31 } \qquad
  \SkewTableau{1,2}{ 45,27,31 }
  \end{example}

  Even though it is possible to draw compositions, and tableaux of
  composition shape, compositions are not supported in the sense that
  many of the options/keys assume that the diagrams are of partition
  shape.  In addition, some options assume that tableaux and diagrams do
  not contain empty rows, and that the diagrams are \textit{connected}.
  (For disconnected tableaux and diagrams, see the
  \keyword{\Multidiagram} and \keyword{\Multitableau} commands in
  \autoref{S:Multitableau}.)

  % --------------------------------------------------------------
  \subsubsection{Drawing order}
  \label{SS:order}
  \index{tableau!drawing order}

  When \TikZ constructs a picture, later objects are drawn over the top
  of earlier ones, which can obscure earlier features. Consequently,
  for more complicated diagrams and tableaux it helps to know the order
  in which the different parts of these diagrams are drawn, which is the
  following:
  \begin{itemize}[nosep]
    \item First, the tableau entries that use the default styling are
    placed, in the order that they appear in the \textit{tableau
    specifications}.
    \item Secondly, the styled tableau entries are placed, in the order
    that they appear in the \textit{tableau specifications}.
    \item Next, the nodes in any \option{ribbons} and ribbon
    \option{paths} are placed, again in the order that the ribbons are
    listed. For each ribbon, first the border of the ribbon is drawn,
    together with any styling for the ribbon, and then the boxes in the
    ribbon, together with any styling, are added.
    \item When enabled by \option{skew border}, the skew border is
    drawn, after which the border of the tableau is drawn.
    \item The \option{dotted rows} and \option{dotted cols} keys are
    applied to replace the specified rows and columns with dots.
    \item Finally, any \option{snobs} are drawn, in the order that they are listed.
  \end{itemize}

  All of the tableau and Young diagram commands use this drawing order.

  % --------------------------------------------------------------
  \subsubsection{Special characters}
  \index{tableau!special characters}
  \index{tableau!comma}
  \index{tableau!*}
  \index{tableau![}

  The characters  \textsf{[}, \textsf{,}, and \textsf{*} have
  special meanings in the \textit{tableau specifications}. Enclose these characters in
  braces when you want to use these characters as entries in a tableau:

  \begin{example}
  \Tableau{ 13{,}5], 8{[}, {*} }
  \end{example}

  There is no need to enclose~\keyword{]} in braces because it only
  becomes special when it has a partner~\keyword{[}.

  % --------------------------------------------------------------
  \subsection{Young diagrams}
  \label{S:Diagram}
  \indexcmd{Diagram}
  \index{young diagram@Young diagram|see{diagram}}
  \index{diagram}
  \index{diagram}
  \index{partition!diagram}

  A \emph{Young diagram}, or simply a \emph{diagram}, is an unlabelled
  tableau. Diagrams can be drawn using the \keyword{\Tableau} command:

  \begin{example}
  \Tableau{~~~,~~~,~~,~}
  \end{example}

  This approach works, but it is cumbersome, hard to proofread, and easy
  to get wrong.  For this reason, \aTableau provides the
  \keyword{\Diagram} command.  The \keyword{\Diagram} command uses
  almost the same syntax as the \keyword{\Tableau} command:

  \keyword{\Diagram (x,y) [options] {partition}}

  As with the \keyword{\Tableau} command, the $(x,y)$-coordinates are
  needed if, and only if, the diagram is inside a \keyword{tikzpicture}
  environment.  The \keyword{\Diagram} command allows the
  \emph{partition} to be specified as either a comma-separated list of
  \emph{weakly decreasing positive integers}, or by using \emph{exponential
  notation}:
  \index{exponential notation}
  \index{partition!exponential notation}
  \index{diagram!exponential notation}

  \begin{example}
  \Diagram{4,3^2,1^3} \qquad
  \Diagram{4,3,3,1}
  \end{example}

  Internally, \keyword{\Diagram} actually does something like the first
  example in this section, so all of the options for the
  \keyword{\Tableau} command can be used with \keyword{\Diagram}. Rather
  than repeating all of the \keyword{\Tableau} options in this section,
  we highlight the more interesting options together with some new,
  diagram specific, options.

  \begin{samepage}
  Just like the \keyword{\Tableau} command, \keyword{\Diagram} supports the
  four different conventions for diagrams, with \option{english} being
  the default:

  \option[-]{english}
  \option*[]{french}
  \option*[]{ukrainian}
  \option*[]{australian}
  \end{samepage}

  \begin{example}
  \Diagram[australian]{3,2,1} \qquad
  \Diagram[australian]{3,1^3}
  \end{example}

  \begin{example}
  % English is the default convention
  \Diagram[english]{3,2,1} \qquad
  \Diagram{3,1^3}
  \end{example}

  \begin{example}
  \Diagram[french]{3,2,1} \qquad
  \Diagram[french]{3,1^3}
  \end{example}

  \begin{example}
  \Diagram[ukrainian]{3,2,1} \qquad
  \Diagram[ukrainian]{3,1^3}
  \end{example}

  \option[true]{border}[true/false]
  \option*[false]{no border}[false/true]

  If \option{border} is true, which it is by default, then the border of
  the tableau is drawn.  If \option{border} is false, then the border of
  the tableau is not drawn. Notice that the border refers only to the
  outside border of the diagram, and not to the internal borders of the
  boxes on the diagram, which can be disabled using
  \option{no boxes}.

  \begin{example}
  \Diagram[border=false]{3,2,1} \qquad
  \Diagram[no border, french]{3^2,1}
  \end{example}

  The \option{no border} option, and similar options, are provided only
  for completeness because it is the inverse of \option{border}. That
  is, \option{no border} is equivalent to \option{border=false}, and
  \option{no border=false} is equivalent to \option{border=true}.

  \option[true]{boxes}[true/false]
  \option*[false]{no boxes}[false/true]

  The \option{boxes} and \option{no boxes} keys control whether or not
  the internal wall around the boxes are drawn, with these two keys
  being inverses of each other. By default, \option{boxes} is true, so
  the internal boxes are drawn.  When \option{no boxes} is in force
  (equivalently, \option{boxes} is false), the internal walls around
  boxes are not drawn.

  \begin{example}
  \Diagram[no boxes]{3,2,1} \qquad
  \Tableau[french, no boxes, no border]{123,45}
  \end{example}

  \option[false]{conjugate}[false/true]

  Draws the conjugate diagram, which has the rows and columns
  interchanged.

  \begin{example}
  \Diagram[conjugate]{3,2,2} \qquad
  \Diagram[ukrainian, conjugate]{2,1^3}
  \end{example}

  \option[]{dotted cols=column indices}[comma separated list of columns]
  \option*[]{dotted rows=row indices}[comma separated list of rows]

  The \option{dotted rows} and \option{dotted cols} keys draw diagrams
  where the specified rows and columns are replaced with dots. This
  makes it possible to draw \emph{generic} diagrams, where the number of
  rows and columns is not fully specified.

  \begin{example}
  \aTabset{scale=0.7}
  \Diagram[dotted rows={2,4,5}]{4^2,2^2,1^3} \qquad
  \Diagram[dotted cols={5,3,4},
    no boxes] {6,5,4,3,2,1}
  \end{example}

  As shown, if row and columns indices are in increasing
  \textit{consecutive} order, then the corresponding rows or columns are
  replaced as a block. Non-consecutive indices are treated separately.

  \begin{example}
  \aTabset{scale=0.7, no boxes}
  \Diagram[dotted rows={4,5,3}] {6,5,5,5,2,1} \qquad
  \Diagram[dotted cols={5,3,4}] {6,5,5,5,2,1}
  \end{example}

  \index{tableau!dotted rows}
  \index{tableau!dotted cols}
  The \option{dotted rows} and \option{dotted cols} keys work by
  overwriting the content on the specified rows and columns, so
  anything that was drawn in these rows and columns will disappear.  The
  \option{dotted rows} and \option{dotted cols} keys can be used with
  the \keyword{\Tableau} command, with the caveat that this will remove
  all boxes, with their contents, in the specified rows and columns.

  When the \option{dotted rows} and \option{dotted cols} keys are
  used together, the \option{dotted rows} are applied first.

  \begin{example}
  \Diagram[scale=0.7, no boxes,
    dotted rows={4,5},
    dotted cols={4,5},
  ] {6,5,5,5,2,1}
  \end{example}

  When using these two keys together, you need to be careful when
  choosing which rows and columns to remove because, otherwise,
  artefacts can remain. Ideally, the endpoints of the rows and columns
  being removed should not overlap. For example, consider the two
  pictures:

  \begin{example}
  \aTabset{scale=0.7, no boxes}
  \Diagram[dotted cols={3}]{6,5,5,5,2,1}\qquad
  \Diagram[dotted cols={4,5,3}, dotted rows={4,5,3},
    tikz after={ \draw[red](0.9,-1.4)circle(0.2); }
  ]{6,5,5,5,2,1}
  \end{example}

  The circled horizontal dots in the right-hand diagram are a ``user
  error'', not a bug. The left-hand diagram shows that these dots come
  from column $3$ with \option{dotted cols} \emph{after} first removing
  rows 4--5 with \option{dotted rows}.

  The \option{dotted rows} and \option{dotted cols} keys are not
  designed to be used with the \option{conjugate} key.

  \option[]{ribbons=ribbon specifications}[list of ribbon specifications]
  \option*[]{paths=path specifications}[list of ribbon specifications]
  \option*[]{snobs=snobs specifications}[list of ribbon specifications]

  Use \option{ribbons}, \option{paths} and \option{snobs} to add
  ribbons, ribbon paths and snobs to a diagram, respectively. The
  \option{snobs} are drawn after the border of the diagram.

  \begin{example}
  \Diagram[french, ribbons={*22*c_4*r*r}]{3,2^2,1} \qquad
  \end{example}

  As described in \autoref{S:Tabloid}, there are additional keys that
  control the behaviour of paths, ribbons and snob. For example, the
  keys \option{path style}, \option{path box} and \option{path box
  style} are associated to \option{paths}.

  \begin{example}
  \Diagram[ukrainian, inner style=dotted,
    styles={b={fill=blue!20,opacity=0.4},
            R={red,text=red,thick},
            B={blue,text=blue,thick},
    }, path style={thick},
    path box=\bullet,
    paths={(R)15rccrc[b]rc,(B)43[b]rcc}
  ]{5^2,3^3}
  \end{example}

  A description of the ribbon specifications, with examples, can be
  found in \autoref{S:RibbonTableaux}, which describes the
  \keyword{\RibbonTableau} command.

  % --------------------------------------------------------------
  \subsubsection{Diagrams with entries}\label{SS:entries}
  \index{contents}
  \index{first tableau}
  \index{hooks}
  \index{last tableau}
  \index{residues}
  \index{first tableau}
  \index{tableau!contents}
  \index{tableau!first}
  \index{tableau!hooks}
  \index{tableau!last}
  \index{tableau!residues}

  The \option{entries} key provides a shortcut for drawing some common
  tableaux of a specified shape.

  \option[]{entries=value}[contents, first, hooks, last, residues]

  The possible choices are:

  \begin{itemize}

    \item\option{entries=first} print the \emph{first}
    \textbf{standard} tableau of this shape, with respect to the Bruhat
    order. This is the tableau that has the numbers $1,2\dots,n$ entered
    in order from top left to right along down successive rows (with
    appropriate modifications for the non-English conventions).

    \index{first tableau}
    \index{tableau!first}
    \index{entries!first}
  \begin{example}
  \Diagram[entries=first]{3,2,2} \qquad
  \Diagram[ukrainian, entries=first]{3^2,2}
  \end{example}

    \item\option{entries=last} print the \emph{last} \textbf{standard}
    tableau of this shape, with respect to the Bruhat order. This is the
    tableau that has the numbers $1,2\dots,n$ entered in order from top
    to bottom down successive columns (with appropriate modifications
    for the non-English conventions).

    \index{last tableau}
    \index{tableau!last}
    \index{entries!last}
  \begin{example}
  \Diagram[entries=last]{3,2,2} \qquad
  \Diagram[australian, entries=last]{2^2,1}
  \end{example}

    \item\option{entries=hooks} print the diagram where each box is
    labelled by the corresponding \emph{hook length}. If $\lambda$ is a
    partition and $\lambda'$ is its conjugate partition then the hook
    length of the box in row~$i$ and column~$j$ is
    $\lambda_i-i+\lambda'_j-j+1$.

    \index{tableau!hook lengths}
    \index{hook lengths}
    \index{entries!hook lengths}
  \begin{example}
  \Diagram[entries=hooks]{3,2,2} \qquad
  \Diagram[ukrainian, entries=hooks]{2^2,1}
  \end{example}

  \item\option{entries=contents} print the diagram where each box is
    labelled by its \emph{content}, which is the row index minus the
    column index.

    \index{tableau!contents}
    \index{contents}
    \index{entries!contents}
  \begin{example}
  \Diagram[entries=contents]{3,2,1} \qquad
  \Diagram[french, entries=contents]{2^2,1}
  \end{example}
  \AddKeyLink{contents}
  \AddKeyLink{residues}
  \AddKeyLink{charge}

  \smallskip\noindent Inside a tableau, the default minus sign for a
  negative number looks too long, so \aTableau uses
  \keyword{\shortminus} for negative contents. For example, it prints
  \keyword{\shortminus 1} instead of \keyword{-1}, which look like
  $\shortminus1$ and $-1$, respectively.\indexcmd{shortminus}

  \item\option{entries=residues} prints the diagram where each box is
  labelled by its \emph{residue}, which is the row index minus the
  column index modulo an integer~$e\ge2$, which must also be supplied.

  \index{residue}
  \index{tableau!residues}
  \index{affine quiver}
  \begin{example}
  \Diagram[entries=residues, e=2]{3,2,1} \qquad
  \Diagram[french, entries=residues, e=3]{3,2,1}
  \end{example}

  \noindent\AddKeyLink{cartan}\AddKeyLink{e}
    By default, residue sequences for affine type $A^{(1)}_e$, or the symmetric
    group and friends, are given. In addition, residues in affine types
    $C^{(1)}_e$, $A^{(2)}_{2e}$ and $D^{(2)}_e$, which can be accessed
    using \option{cartan=C}, \option{cartan=AA}, and \option{cartan=DD},
    respectively.

    \smallskip
  \begin{example}
  % affine type C
  \Diagram[french, entries=residues,
    e=2, cartan=C] {8,4,2}
  \end{example}

  \begin{example}
  % twisted affine type A
  \Diagram[ukrainian, entries=residues,
     e=2, cartan=AA] {8,4,2}
  \end{example}

  \begin{example}
  % twisted affine type D
  \Diagram[australian, entries=residues,
    e=2, cartan=DD] {8,4,2}
  \end{example}

  \end{itemize}
  Use \option{charge} to add offsets to the \option{contents} and \option{residues}.

  \begin{example}
  \Diagram[entries=contents, charge=2]{3,2,1} \qquad
  \Diagram[entries=residues, e=3, charge=1]{3,2,1}
  \end{example}

  % --------------------------------------------------------------
  \subsection{Skew tableaux and skew diagrams}
  \label{S:Skew}
  \indexcmd{SkewTableau}
  \indexcmd{SkewDiagram}
  \index{diagram!skew}
  \index{partition!skew}
  \index{tableau!skew}

  A partition $\mu$ is \emph{contained} in another partition $\lambda$, which
  is written as~$\mu\subseteq\lambda$, if $\mu_k\le\lambda_k$, for
  $k\ge0$. If $\mu\subseteq\lambda$ then the \emph{skew partition}
  $\lambda/\mu=\set{(r,c)\in\lambda|(r,c)\notin\mu}$ is the set of nodes
  that are in~$\lambda$ and not in~$\mu$. In this manual, $\mu$ is
  the \emph{inner shape} and $\lambda/\mu$ is the outer shape. A
  \emph{skew tableau} is a labelling of the nodes in the diagram of a
  skew partition. Skew partitions and skew tableau can be drawn using
  the commands:

  \keyword{\SkewDiagram (x,y) [options] {inner shape} {outer shape}}\newline
  \keyword{\SkewTableau (x,y) [options] {inner shape} {skew tableau specifications}}

  \index{exponential notation}
  As with the previous commands, the $(x,y)$--coordinates should be given
  if, and only if, the picture is inside a \keyword{tikzpicture}
  environment.  The inner and outer shapes can either be given as a
  comma-separated list of non-negative integers, or in exponential notation.

  Skew tableau and skew diagrams should always be connected. Use
  \keyword{\Multidiagram} and \keyword{\Multitableau} to draw
  disconnected diagrams.

  In fact, the \keyword{\SkewTableau} is just an alias for the
  \keyword{\Tableau}, using the \option{skew} key to set the inner
  shape. For this reason, all of the options for the \keyword{\Tableau}
  command can be used with \keyword{\SkewTableau}. The entries in a
  \keyword{\SkewTableau} are specified in exactly the same way as in the
  \keyword{\Tableau} command. In particular, the entries of skew
  tableaux can, optionally, be given style prefixes, both using a~$*$ to
  add the current \option{star style}, or arbitrary \TikZ-styling using
  \keyword{[...]}.

  \begin{example}
  \tikzset{R/.style={fill=red!20,text=red}}
  \SkewTableau[french]{1^2}{13,5*7,[R]79} \qquad
  \Tableau[french, skew={1^2}]{13,5*7,[R]79}
  \end{example}

  Similarly, the \keyword{\SkewDiagram} command is an alias for the
  \keyword{\Diagram} command, with a \option{skew} key.

  \begin{example}
  \SkewDiagram[french]{1^2}{2^3} \qquad
  \Diagram[french, skew={1^2}]{2^3}
  \end{example}

  As the \keyword{\SkewTableau} and \keyword{\SkewTableaua} commands
  are shortcuts, all of the options for the \keyword{\Tableau} and
  \keyword{\Diagram} commands can be used with the
  \keyword{\SkewTableau} and \keyword{\SkewDiagram} commands,
  respectively. In addition, the following options are supported:

  \option[false]{skew border}[true/false]
  \option*[true]{no skew border}[false/true]

  These two keys are inverse to each other.  When
  \option{skew border} is \textsf{true}, the border of the (inner) skew
  shape is drawn.

  \begin{example}
  \tikzset{R/.style={fill=red!20,circle}}
  \SkewTableau[skew border]{2,1^2}{12[R]3,45,67}
  \SkewTableau[french, skew border]{2,1^2}{12[R]3,45,67}
  \end{example}

  Note that \option{skew border} only draws the border of the skew shape
  and not the boxes inside the inner skew shape. Use \option{skew boxes}
  to draw the interior walls of the boxes in the skew shape.

  \option[draw=\aTableauColour{Skew}, fill=\aTableauColour{SkewFill}, thick]{skew border style=style}[\TikZ-styling]

  Use \option{skew border style} to change the style of the skew border.
  Since the skew border is only drawn when \option{skew border} is set,
  the \option{skew border style} does not have any effect unless
  \option{skew border} has been set to \textsf{true}.

  \begin{example}
  \SkewTableau[skew border style={draw=yellow},
      skew border]{2,1^2}{123,45,67} \qquad
  \SkewDiagram[skew border style={dashed,fill=red!10},
      skew border]{1^2}{2^3}
  \end{example}

  \option[false]{skew boxes}[false/true]
  \option*[true]{no skew boxes}[true/false]
  \option*[thin, fill=\aTableauColour{SkewFill}]{skew box style=style}[\TikZ-styling]

  The \option{skew boxes} key adds walls to the boxes inside the inner
  partition of a skew shape. Use \option{skew box style} to change the
  default shading of these boxes.

  \begin{example}
  \SkewTableau[skew boxes] {2,1^2}{123,45,67}
  \SkewTableau[skew boxes,
    skew box style={dash pattern=on 1pt off 2pt}
  ] {2,1^2}{123,45,67,89}
  \end{example}

  Using the options \option{skew border} and the \option{skew boxes}
  together gives the skew (inner) shape both inner and outer borders.

  \begin{example}
  \SkewTableau[skew border, skew boxes] {2^2,1}{123,45,67}
  \end{example}

  % --------------------------------------------------------------
  \subsection{Shifted tableaux and shifted diagrams}
  \label{S:Shifted}
  \indexcmd{ShiftedTableau}
  \indexcmd{ShiftedDiagram}
  \index{diagram!shifted}
  \index{strict partition}
  \index{partition!strict}
  \index{shifted tableau}
  \index{tableau!shifted }

  \emph{Strict partitions}, which have strictly increasing parts, and
  \emph{shifted tableaux}, which are of strict partition shape, appear
  in several places in representation theory, such as in the study of
  projective representations of the symmetric groups. These tableaux and
  diagrams are drawn with row~$r$ shifted by $(r-1)$-units along the
  row, so that the first box in each row has content~$0$. These diagrams
  and tableaux can be drawn using the following commands:

  \keyword{\ShiftedDiagram (x,y) [options] {partition }}\newline
  \keyword{\ShiftedTableau (x,y) [options] {tableau specification}}

  As usual, the $(x,y)$-coordinates are necessary if, and only if, these
  commands are used inside a \keyword{tikzpicture} environment. The
  partition in a~\keyword{\ShiftedDiagram} can be given using
  exponential notation and the tableau specification for a
  \keyword{\ShiftedTableau} can include the usual optional style
  prefixes.

  \begin{example}
  \tikzset{R/.style={fill=red!20,circle}}
  \ShiftedTableau{12[R]3,45,*6}  \qquad
  \ShiftedDiagram[french]{3,2,1}
  \end{example}

  In the literature,  shifted tableaux and shifted diagrams almost
  always have strict partition shape, however, this is not enforced by
  these commands. Under the hood, the \keyword{\ShiftedTableau} is the
  \keyword{\Tableau} command with the \option{shifted} option set
  to true. Similarly, \keyword{\ShiftedDiagram} is the same as using
  the \keyword{\Diagram} command with the \option{shifted} option. Consequently,
  all of the options for the \keyword{\Tableau} and \keyword{\Diagram}
  commands can be used with \keyword{\ShiftedTableau} and
  \keyword{\ShiftedDiagram}, respectively.  In particular, options like
  \option{skew boxes} can be used to highlight the shifted part of the
  diagram

  \begin{example}
  \ShiftedTableau[skew boxes]{1*23,4*5} \quad
  \ShiftedDiagram[skew border]{3,2}
  \end{example}

  The \option{entries} key works as expected for shifted diagrams:

  \begin{Example}
  \ShiftedDiagram[entries=contents]{3,2,1} \qquad
  \ShiftedDiagram[entries=first]{4,2,1} \qquad
  \ShiftedDiagram[entries=hooks]{4,2,1} \qquad
  \ShiftedDiagram[entries=last]{4,2,1} \qquad
  \ShiftedDiagram[entries=residues, e=3]{4,2,1}
  \end{Example}

  The definition of hook lengths for shifted diagrams can be found, for
  example, in \cite{Konvalinka:ShiftedHookLengths}.

  % --------------------------------------------------------------
  \subsection{Tabloids}
  \label{S:Tabloid}
  \indexcmd{Tabloid}

  A \emph{tabloid} is an equivalence class of tableau, where two tableaux
  are equivalent if they have the same entries in each row, up to a
  permutation. In the literature, tabloids are usually drawn with lines
  above and below each row, and without side borders. Tabloids can be
  drawn with the \keyword{\Tabloid} command.

  \keyword{\Tabloid (x,y) [options] {partition }}

  The \keyword{\Tabloid} command functions in the same way as the
  \keyword{\Tableau} command, except that it draws tabloids. In fact,
  the \keyword{\Tabloid} is a special case of the \keyword{\Tableau}
  command with the \option{tabloid} option set.

  \begin{example}
  \tikzset{R/.style={fill=red!20,circle,minimum size=5mm}}
  \Tabloid[ukrainian]{12[R]3,45}\quad
  \Tabloid[english]  {12[R]3,45}\quad
  \Tabloid[french]   {12[R]3,45}
  \end{example}

  The \aTableau package does not provide a dedicated command for tabloid
  diagrams, however, they can be drawn using the \keyword{\Diagram}
  command together with the \option{tabloid} option:

  \begin{example}
  \Diagram[tabloid]{3^2,1} \qquad
  \Diagram[tabloid, french]{3,1^2}
  \end{example}

  Some authors work with \emph{column tabloids}. Such tableaux do not
  have a dedicated \aTableau command, however, they can be drawn using the
  \option{conjugate} option:  \index{tabloids!column}

  \begin{example}
  \Tabloid[conjugate]{123,4} \qquad
  \Diagram[conjugate, tabloid]{3,1} \qquad
  \end{example}

  Similarly, skew tabloids and shifted tabloids can be drawn using the \option{skew}
  option.
  \label{tabloids!skew}

  \begin{example}
  \Tabloid[skew={2,1}]{19{11},29{10},6} \qquad
  \Tabloid[conjugate, skew={0,1,2}]{19{11}7,29{10},6}
  \end{example}


  % --------------------------------------------------------------
  \subsection{Ribbon tableaux}\label{S:RibbonTableaux}
  \indexcmd{RibbonTableau}
  \index{ribbon!head}
  \index{tableau!ribbon tableau}

  A \emph{ribbon} in a tableau, or a diagram, is a connected strip of
  boxes~$R$ that are totally ordered by their contents. (Recall from
  \autoref{SS:entries} that the content of the box in row~$a$ and
  column~$b$ is $b-a$.) Therefore, a ribbon does not contain a
  \Diagram[scale=0.3, border=false, align=bottom]{2^2}-square, and if
  $(a,b)\in R$ then at most one of $(a+1,b)$ and $(a,b-1)$ belongs
  to~$R$.  The \emph{head} of a ribbon~$R$ is the unique node of maximal
  content.  A \emph{ribbon tableau} is a tableau that is tiled by
  ribbons. A box is a ribbon of length~$1$, so every tableau is a ribbon
  tableau.

  \keyword{\RibbonTableau (x,y) [options] {ribbon specifications}}

  Like the other commands, the $(x,y)$-coordinates are optional and are
  required if, and only if, the diagram is being used as part of a
  \keyword{tikzpicture} environment. All options for the
  \keyword{\Tableau} command can be used for the ribbon command, so we
  refer to \autoref{S:Tableau} for a description of the available
  options.

  % --------------------------------------------------------------
  \subsubsection{Ribbon specifications}
  Unlike the \keyword{\Tableau} command, the entries of a ribbon tableau
  are given by ribbon specifications (rather than tableau)
  specifications. Ribbon specifications are also used by the three keys
  \option{ribbons}, \option{paths}, and \option{snobs}, which were
  introduced in \autoref{SS:TableauOptions}.

  To understand the ribbon specifications, observe that if $(a,b)$ is
  the head of a ribbon, then we can walk along the ribbon by specifying
  whether the \textsf{r}ow index increases or the \textsf{c}olumn index
  decreases. That is, a ribbon is uniquely determined by specifying its
  head $(a,b)$ together with a sequence of $\textsf{r}$'s and
  \textsf{c}'s to indicate when the row index increases, or the column
  index decreases, respectively. For example, vertical and horizontal
  ribbons are given by a sequences of \textsf{r}'s and \textsf{c}'s,
  respectively:

  \begin{example}
  \aTabset{align=top, no border}
  \RibbonTableau{11rrrr} \qquad
  \RibbonTableau{16ccccc}
  \end{example}

  The border of a ribbon tableau is the smallest (skew) partition that
  contains all of its ribbons.  In the last example shows, the
  \option{no border} key is used to disable the ribbon tableau border so
  that we can better see the ribbon, which is the full diagram in these
  two cases.

  The following diagram shows a see-sawing ribbon with head in row~$1$
  and column~10, and tail in row~$5$ and column~$1$.

  \begin{example}
  \RibbonTableau{1{10}crcrcrcrccccc}
  \end{example}

  As with tableau entries, each box in a ribbon can be preceded with a
  style specification, either using \keyword{*} for the \option{star
  style} or using \keyword{[...]} for more customised style
  specifications.

  \begin{example}
  \RibbonTableau{
      *1{10}c*rc*rc*rc*rc*cc*cc,
      *18c[fill=red]rc[fill=blue]rc*rc*cc[fill=teal]c,
      [fill=green]16c*rc*rc*cc,
      *14c[{draw=purple,ultra thick}]cc[fill=orange]r,
      [{fill=yellow,draw=red,ultra thick}]23
  }
  \end{example}

  In particular, the yellow box in this example highlights that boxes
  are ribbons of length~$1$. Most of the time, you want to style the
  entire ribbon rather than styling the individual boxes in the ribbon.
  Optionally, the style for the entire ribbon can be given inside
  parentheses $(...)$ at the \textit{start} of the ribbon. Unlike the
  other style specifications, the $*$-shorthand for the \option{star
  style} cannot be applied to an entire ribbon.

  \begin{example}
  \RibbonTableau{
      (fill=red)1{10}crcrcrcrccc_Ccc,
      (fill=orange)18crcrc*rcc_Bcc,
      (fill=yellow)16crcr[fill=cyan]c_Acc,
      (fill=green)14cccr,
      (fill=blue)23
  }
  \end{example}

  Any styles for a ribbon are applied to the region bounded by  the
  closed path given by the border of the ribbon. The ribbon style is
  applied first, after which any styled boxes in ribbon are drawn.  As
  this example shows, the optional style of any box in the ribbon takes
  precedence over the style of the full ribbon.

  Finally, it is often necessary to label some of the boxes in the
  ribbon. Optionally, text can be added to the boxes in a ribbon by
  supplying it as a \textit{subscript} to the corresponding entry in the
  ribbon specification:

  \begin{example}
  \RibbonTableau[french]    { 12_1 c_2 r_3 } \quad
  \RibbonTableau[ukrainian] { 12_a c_b r_c } \quad
  \RibbonTableau[australian]{ 12_{1_a} c_{2_b} r_{3_c} }
  \end{example}

  Here we have put spaces between the entries for the different boxes
  in the ribbon for clarity. There are often many ways to draw the same
  ribbon tableau, especially if the entries have styling.

  \begin{example}
  \tikzset{Fill/.style={fill=#1!20!white},
    B/.style={Fill=blue}, R/.style={Fill=red}, Y/.style={Fill=yellow}}
  \RibbonTableau{ [B]12_1  [R]c_2   [Y]r_3} \qquad
  \RibbonTableau{ [B]12_1, [R]11_2, [Y]21_3 } \qquad
  \Tableau{       [B]2     [R]1,    [Y]3 }
  \end{example}

  As shown, the easiest way to draw this particular tableau is using the \keyword{\Tableau} command.

  There is no dedicated syntax for putting a label on the boundary
  between two boxes in the diagram, but this can be done using the named
  anchors of~\autoref{E:NodeNames}.

  \begin{example}
  \RibbonTableau[shifted, tikz after={
      \node at (A-1-2.east){$1$};
      \node at (A-2-2.east){$2$};
  }]{ 14c, 12c, 23c }
  \end{example}

  Here is an example of a $3$-ribbon skew tableau.

  \begin{example}
  \RibbonTableau[skew={2^2,1^2}, skew border]{
     14c_0r,15r_1c, 18c_7c,26r_9r,
     34c_2c,35r_8r,42r_3c, 44c_4r,
     54r_6c,62c_5r
  }
  \end{example}

  We have now seen the full ribbon specifications. To summarise:
  \begin{itemize}[nosep]
    \item The ribbon specification starts with optional \TikZ-styling
    for the ribbon (or path), enclosed in $(...)$.
    \item The first entry in the ribbon is given as
    \textsf{a}\textsf{b}, if the head of the ribbon is in row
    \textsf{a} and column \textsf{b}. The remaining entries in the
    ribbon are specified by a sequence of \textsf{r}'s and \textsf{c}'s,
    depending on whether the next box in the ribbon is given by
    increasing the row index, or decreasing the column index,
    respectively.
    \item Optionally, each box in the ribbon, including the head, can be
    given \TikZ-styling by prefixing it with a $*$, which gives the
    box the \option{star style}, or by giving \TikZ-styling keys inside
    square brackets $[...]$.
  \end{itemize}

  Finally, when drawing the border, all ribbons in a
  \keyword{\RibbonTableau} are assumed to be contained in the smallest
  Young diagram that contains the ribbons. This said, it is up to you to
  ensure that the boxes in the ribbon do not have negative column
  indices.  Inside a \keyword{\Diagram} or \keyword{\Tableau} command,
  any ribbons added using \option{paths}, \option{ribbons}, or
  \option{snobs} do not affect the border of the diagram or tableau. You
  can draw paths, ribbons, and snobs so that their boxes are either
  inside or outside the border of the diagram or tableau. It can be
  useful to place boxes outside of the diagrams.

  \option[]{paths=path specifications}[comma separated ribbon specifications]<ribbons>
  \option*[]{ribbons=ribbon specifications}[comma separated ribbon specifications]<ribbons>
  \option*[]{snobs=snob specifications}[comma separated ribbon specifications]<ribbons>

  You can add \option{ribbons}<ribbons>, \option{paths}<ribbons> and
  \option{snobs}<ribbons> to a \keyword{\RibbonTableau} in exactly the
  same way that you add them to \keyword{\Tableau}. At first sight, it
  seems unnecessary to add \option{ribbons}<ribbons> and
  \option{snobs}<ribbons> because ribbon tableaux are composed of
  ribbons. The point is that any ribbon added to a tableau using
  \option{ribbons}<ribbons> or \option{snobs}<ribbons> is not assumed to
  be inside the border of the tableau.

  \begin{example}
  \RibbonTableau[
    styles={A={fill=green!10, text=teal}, R={fill=yellow!10, text=red}},
    ribbons={[A]15_A,[A]24_A,[A]41_A, [R]14_R,[R]33_R}
  ]{11, (fill=blue!10)12rc, 14crrcc}
  \end{example}

  Here is a cute use of \option{paths} to add arrows to a skew
  $2$-ribbon tableau:

  \begin{example}[lefthand width=0.35\textwidth]
  \RibbonTableau[skew={4,3,2,1},
    styles={L={red,semithick,->}, R={red,semithick,<-}},
    paths ={ (L)15c, (L)17c, (L)24c, (L)26c,
             (R)32r,(R)41r,(R)52r,(R)61r, (R)34c},
    tikz after = { % node labels: row/col/label
      \foreach \r/\c/\l in {1/6/0,1/8/6,2/5/3,2/7/8}
        { \node at (A-\r-\c.west){\l};}
      \foreach \r/\c/\l in {5/1/1,4/2/2,3/3/5,7/1/4,6/2/7}
        { \node at (A-\r-\c.south){\l};}
    }
  ]{ 16c,18c, 25c, 27c, 33r, 42r, 51r, 62r,71r }
  \end{example}

  The arrows are added by the \option{paths} key using the two styles
  \textsf{L} and \textsf{R}, which draws ``left'' and ``right'' arrows
  between the two boxes, where left and right is with respect to the
  line-order of the two boxes in the corresponding draw command.  The
  two \keyword{\foreach}-statements might look complicated, but they are
  just loops that add the entries to the tableau. These loops set
  \keyword{\r}, \keyword{\c} and \keyword{\l} equal to the row index,
  column index and label, respectively, using the two lists of
  \keyword{\r}\!/\!\keyword{\c}\!/\!\keyword{\l} triples given above.

  \option[draw=\aTableauColour{Inner}, thin]{ribbon style=style}{\TikZ-styling}
  \option*[]{ribbon box style=style}{\TikZ-styling}

  Use the \option{ribbon style} and \option{ribbon box} keys to
  append to the default styling of the ribbons, and the boxes in the
  ribbons, respectively. By default, the full ribbon has inner wall
  styling, but the boxes in ribbons have no styling. If you want inner
  walls, then they can be added using the \option{ribbon box} key.

  \begin{example}
  \RibbonTableau[ribbon style={draw=red, thick,
    fill=orange!20}]{16_1r_2c_3c_4c_5c_6r_7c_8}
  \RibbonTableau[ribbon box style={fill=orange!20,
    draw=magenta}]{16_1r_2c_3c_4c_5c_6r_7c_8}
  \end{example}

  Ribbons are drawn by first drawing the full ribbon, applying
  \option{ribbon style}, then the empty boxes in the ribbon are drawn,
  and, finally, any boxes in the ribbon that have styling or an entry
  are added. In particular, if the style \option{ribbon box} draws an
  outline around a box then this has precedence over any outline added
  by \option{ribbon style} since \option{ribbon box} is applied last.

  \option[]{ribbon box=text}{default box contents}

  You can add content to the boxes in a ribbon using subscripts in
  the ribbon specifications. Alternatively, you can use
  \option{ribbon box} if you want to add the same
  content to each box in the ribbon.

  \begin{example}
  \RibbonTableau[ribbon box=\times]{14rccrc}\qquad
  \RibbonTableau[ribbon box=\bullet]{14_1rc_2cr_3c}
  \end{example}

  As shown, whenever a box entry is given as a subscript it is used
  instead of the value of \option{ribbon box}. If a ribbon
  tableau contains multiple ribbons, then \option{ribbon box} is used
  for all ribbons.

  Similarly, there are keys \option{path style}, \option{path box},
  \option{path box style}, \option{snob style}, \option{snob box} and
  \option{snob box style} that work in exactly the same way. Most of the
  \keyword{\Tableau} keys work as expected for ribbon tableaux.

  % --------------------------------------------------------------
  \subsection{Multidiagrams and multitableaux}\label{S:Multitableau}
  \indexcmd{Multidiagram}
  \indexcmd{Multitableau}
  \index{diagram!multidiagram}
  \index{tableau!multitableau}

  Multitableaux and multidiagrams are $\ell$-tuples of tableaux and
  diagrams, respectively. Such diagrams appear in many places in representation
  theory, such as when considering representations of wreath products
  or cyclotomic Hecke algebras. For convenience, \aTableau provides
  the \keyword{\Multidiagram} and \keyword{\Multitableau} commands for
  drawing these diagrams, even though it is possible to construct
  such tableaux and diagrams using the \keyword{\Tableau} and
  \keyword{\Diagram} commands.

  \keyword{\Multidiagram (x,y) [options] {multipartition specifications} }\newline
  \keyword{\Multitableau (x,y) [options] {multitableau specifications}   }

  \index{components}
  \index{multidiagram!components}
  \index{multitableau!components}
  Our preferred notation for a multipartition $\blam$ is to write
  $\blam=(\lambda^{(1)}|\lambda^{(2)}|\dots|\lambda^{(\ell)})$, where
  $\ell\ge1$ is the \emph{level} and the partitions
  $\lambda^{(1)},\dots,\lambda^{(\ell)}$ are the \emph{components}
  of~$\blam$. Similarly, a multitableau is written as
  $\T=(\T^{(1)}|\dots|\T^{(\ell)})$, where the tableaux
  $\T^{(1)},\dots,\T^{(\ell)}$ are the \emph{components} of~$\T$.
  Accordingly, in the \keyword{\Multidiagram} and
  \keyword{\Multitableau} commands, the pipe symbol \keyword{|} is used
  to separate components:

  \begin{example}
  \aTabset{scale=0.7}
  \Multidiagram{3,2^2|2,1,1|1}
  \Multitableau[box font=\tiny]
    { 123,45,67 | 89,{10},{11} | {12}{13}{14} }
\end{example}

  As shown, the \keyword{\Multidiagram} command accepts exponential
  notation for multipartitions. The entries in a \keyword{\Multitableau} can
  be given optional styling specifications, exactly as inside the
  \keyword{\Tableau} command.

  \begin{example}
  \tikzset{R/.style={fill=red!10,text=magenta}}
  \Multitableau{1*24,*68 | [R]9{10},{11}}
  \end{example}

  Multitableaux and multidiagrams can, in principle, have arbitrary
  level. In practice, the level is constrained by the page width.

  \begin{example}
  \Multitableau[french]{1,2 | 3,4 | 5,6 | 7,8 | 9,{10}}
  \end{example}

  Both the \keyword{\Multitableau} and \keyword{\Multidiagram} commands
  accept a special component, \keyword{...}, that is used to add dots between the
  separators:

  \begin{example}
  \aTabset{scale=0.7}
  \Multidiagram{3,2^2|...|1}
  \Multitableau[box font=\tiny]
    {123,45,67 | ... | {12}{13}{14}}
  \end{example}

  The \keyword{\Multitableau} and \keyword{\Multidiagram} commands
  accept most of the options of the \keyword{\Tableau} and
  \keyword{\Diagram} commands, together with a few options that are
  specific to multitableaux and multidiagrams. This section describes
  the new options, and highlights how some of the other options are used.

  \option[-]{english}
  \option*[]{french}
  \option*[]{ukrainian}
  \option*[]{australian}

  All four tableaux conventions are supported for multitableaux and
  multidiagrams, with the caveat that all component diagrams use the
  same convention.

  \begin{example}
  \aTabset{scale=0.7}
  \Multidiagram[french]{3,2^2|...|1}
  \Multitableau[ukrainian, box font=\tiny]
    {123,45,67 | ... | {12}{13}{14}}
  \end{example}

  \option[]{conjugate}
  \index{conjugate!multidiagram}
  \index{conjugate!multitableau}
  \index{multidiagram!conjugate}
  \index{multitableau!conjugate}

  The \option{conjugate} option prints the conjugate multitableaux and
  multidiagrams. Conjugation for multitableaux and multidiagrams reverses
  the order of the components and then conjugates the component
  diagrams.

  \begin{example}
  \aTabset{scale=0.6, box font=\scriptsize}
  \Multitableau[conjugate]{ 123,45 | 67,89} \quad
  \Multidiagram[conjugate]{ 3,2 | 1^2 }
  \end{example}

  \option[()]{delimiters}[pair of \LaTeX{} delimiters]
  \option*[(]{left delimiter}[a \LaTeX{} delimiter]
  \option*[)]{right delimiter}[a \LaTeX{} delimiter]

  The delimiter options control the delimiters that are put on the left
  and right of multitableaux and multidiagrams. These options accept any
  \LaTeX{} delimiter.

  \begin{example}
  \aTabset{scale=0.7}
  \Multidiagram[left delimiter=\langle,
      right delimiter={]}]{2,1|1} \quad
  \Multidiagram[delimiters={\|}{\}} ]{1^2|2}
  \end{example}

  To remove a delimiter, set \option{left delimiter} or
  \option{right delimiter} to nothing.

  \begin{example}
  \Multitableau[left delimiter=, right delimiter=]
    {123,45,67|{12}{14}}
  \end{example}

  \option[\textbackslash textendash]{empty=text}[any \LaTeX{}]

  The \option{empty} option determines what is printed when a component
  diagram is empty. By default, the en-dash symbol \textendash{} is
  used.

  \begin{example}
  \Multitableau{ 123,45 | | 67,89} \par\bigskip
  \Multidiagram[empty=$\emptyset$]{ 3,2 | | 2^2 }
  \end{example}

  \begin{warning}
    By default, the \option{empty} key is typeset as text and \emph{not}
    as mathematics, which is the default for the tableau boxes.
  \end{warning}

  \option[]{entries=value}[contents, first, hooks, last, residues]<multitableau>
  \option*[]{charge=offsets}[$|$-separated list of charges]<multitableau>

  \index{contents}
  \index{residues}
  \index{contents!multitableau}
  \index{first tableau!multitableau}
  \index{hooks!multitableau}
  \index{last tableau!multitableau}
  \index{residues!multitableau}
  \index{entries!multitableau}
  \index{multidiagram!show}
  \index{multitableau!contents}
  \index{multitableau!first}
  \index{multitableau!hooks}
  \index{multitableau!last}
  \index{multitableau!residue}

  The \option{entries} key automatically fills in particular types of
  entries for you to draw the requested tableau of the specified shape.
  Use \option{charge} to provide offsets in each component when using
  \option{entries=contents} and \option{entries=residues} {(}the
  \option{charge} is ignored when using \option{entries=first} and
  \option{entries=last}).

  \nopagebreak[0]
  The possible choices are:
  \begin{itemize}

    \item\option{entries=first} print the \emph{first} \textbf{standard} tableau of this
    shape, which has the numbers $1,2\dots,n$ entered in order from top
    left to right along down successive rows (with appropriate modifications for
    the non-English conventions).

  \begin{example}
  \Multidiagram[entries=first]{3,1^2|2^2,1}
  \end{example}

    \item\option{entries=last} print the \emph{last} \textbf{standard} tableau of this
    shape, which has the numbers $1,2\dots,n$ entered in order from top
    to bottom down successive columns, and right to left along components (with appropriate modifications for
    the non-English conventions).

  \begin{example}
  \Multidiagram[entries=last]{3,2,2|2^2}
  \end{example}

    \item\option{entries=hooks} print the diagram where each box in each
    component is labelled by the corresponding \emph{hook length}.

  \begin{example}
  \Multidiagram[entries=hooks]{3,1^2|3^3}
  \end{example}

    \item\option{entries=contents} print the diagram where each box is
    labelled by its \emph{content}, which is the row index minus the
    column index.

  \begin{example}
  \Multidiagram[entries=contents,
    charge={0,2}] {3,2,1|2^2}
  \end{example}

    As shown, the \option{charge} offsets the contents in each
    component.

    \item\option{entries=residues} print the diagram where each box is
    labelled by its \emph{residue}, which is the row index minus the
    column index modulo some integer~$e$, which must also be supplied.

    \index{residues}
    \index{tableau!residues}
  \begin{example}
  \Multidiagram[entries=residues, e=2]{3,2,1|1^3}
  \end{example}

  Specifying the \option{charge} offsets the residues in each component.

  \begin{example}
    \Multidiagram[french, entries=residues, e=3,
      charge={1,2}]{3,2,1|1^2}
  \end{example}

  \medskip\noindent
  As in \autoref{SS:entries}, residues for the other affine Cartan types
  can be specified using \option{cartan=value}.

  \end{itemize}

  \option[]{label=labels}[a $|$-separated list of labels]

  Use \option{label} to add labels to the component diagrams.

  \begin{example}
    \Multitableau[label={1|2}]{123,45,67|{12}{14}}
  \end{example}

  No extra space is allowed for the labels, so use the
  \option{separation} key, or \option{xoffsets} and \option{yoffsets},
  to adjust for wide labels.

  \option[A]{name=text}[character]<multi>

  The \option{name} option is used to change the default prefix for the
  box node names. In a tableau, or diagram, the box names depend on the
  \emph{row} and \emph{column} indices. In a multitableau, or
  multidiagram, the box names depend on the \emph{component}, \emph{row}
  and \emph{column} indices. Explicitly, by default, in multitableaux
  and multidiagrams the names of the take the form \textsf{A-k-r-c},
  where~\textsf{k} is the component index,~\textsf{r} is the row index,
  and~\textsf{c} is the column index of the box.
  \footnote{This follows the convention for nodes in multidiagrams
  from~\cite{EvseevMathas:DeformedKLR}, rather than the somewhat misguided
  convention that the author introduced in~\cite{DJM:cyc}.}
  The \textsf{B} can be changed like using the \option{name} key.

  \begin{example}
  \Multitableau[name=A, tikz after = {
    \draw[red, thick,<->](A-1-2-2.south east) to
         [out=315,in=225] (A-2-2-1.south west);
    }] { 123,4*5 | 67,*89}
  \end{example}

  \option[]{paths=path specifications}[a $|$-separated list of ribbon path specifications]<multi>
  \option*[]{ribbons=ribbon specifications}[a $|$-separated list of ribbon specifications]<multi>
  \option*[]{snobs=snob specifications}[a $|$-separated list of ribbon specifications]<multi>

  The \option{paths}<multi>, \option{ribbons}<multi>,  and
  \option{snobs}<multi> keys work in exactly the same way for
  multitableaux and multidiagrams as they do for tableaux and diagrams,
  except that the ribbons for different components are separated by a
  pipe $|$.

  \begin{example}

  \Multitableau[ribbons={(R)11*r | (Y)12rc},
    styles={R={fill=red!10, opacity=0.5},
            Y={fill=yellow!20, opacity=0.5}
  }] { 123,45 | 67,89}
  \end{example}

  Notice that the $4$ is the first component has disappeared because
  the second node in the ribbon \textsf{(R)11*r} has $*$-styling, which
  writes over the $4$. (One way to recover the missing entry is to use the
  ribbon specification \textsf{(R)11*r\textunderscore4}.)

  \option[]{rows=number}[a decimal number]

  The \option{rows} option sets the number of rows covered by the
  delimiters in multitableaux and multidiagram.  By default, the
  \keyword{\Multitableau} and \keyword{\Multidiagram} set
  the size of the delimiters so that they match the component diagrams, which may
  not be what you want especially when using the \option{australian} and
  \option{ukrainian} conventions. The value of \option{rows} can be a
  decimal number.

  \begin{example}
  \Multidiagram[ukrainian, scale=0.5]{3,2^2|...|2,1^2}
  \end{example}

  \begin{example}
  \Multitableau[rows=3, scale=0.8, box font=\tiny]
       {{1_1}{1_1}{1_1}{2_1},{2_1}{2_1}{2_1} |
        {1_2}{3_3},{2_2} |
        {1_3}{1_3},{2_3}}
  \end{example}

  \begin{example}
  \Multitableau[ukrainian, rows=2.5, scale=0.8, box font=\tiny]
       {{1_1}{1_1}{1_1}{2_1},{2_1}{2_1}{2_1} |
        {1_2}{3_3},{2_2} |
        {1_3}{1_3},{2_3}}
  \end{example}

  \option[true]{separators}[true/false]
  \option*[false]{no separators}[false/true]

  When \option{separators} is true, delimiters are drawn around the
  component diagrams in a multitableaux and a multidiagram and
  separators are drawn between the component shapes, and
  \option{no separators} is the inverse option. By default, separators
  and delimiters are drawn.

  \begin{example}
  \Multitableau[separators, scale=0.8]
    {123,45,67 | 89,{10}{11}}  \par\bigskip
  \Multidiagram[no separators, scale=0.7]
    {3,2^2|2,1^2} \par\bigskip
  \end{example}

  \option[$|$]{separator=text}[character]

  The \option{separator} determines what is printed between the
  component diagrams in multitableaux and multipartition. By default,
  \option{separator=$|$} draws a straight line between the
  components, but any character can be used.

  \begin{example}
  \Multitableau[separator={,}]{ 123,45 | 67,89} \par\bigskip
  \Multidiagram[separator=:]{ 3,2 | 2^2 }
  \end{example}

  \option[\aTableauColour{Main}]{separator colour=colour}[a \LaTeX{} colour]

  The \option{separator colour} sets the colour of the delimiters and
  the separator between component diagrams. By default, the separators
  are the same colour as the diagram border.

  \begin{example}
  \Multitableau[separator colour=red]{123,45 | 67,89} \par\bigskip
  \Multidiagram[separator color=brown, separator=:]
    {3,2 | 2^2 }
  \end{example}

  \option[0.3]{separation=number}[a decimal number (the distance in \unit{cm})]

  The \option{separation} key sets the distance between the separators
  and the component tableaux and diagrams.

  \begin{example}
  \Multitableau[separation=0.1]{123,45 | 67,89} \par\bigskip
  \Multidiagram[separation=0.8]{3,2 | 2^2}
  \end{example}

  \option[]{skew=partitions}{a $|$-separated list of partitions}

  Skew multishapes can be drawn using the \option{skew} option to
  specify a $|$-separated list of inner shapes.

  \begin{example}
    \Multitableau[skew={2,1|1}]{ 123,45 | 7{10},89} \par\bigskip
    \Multidiagram[skew={1|2,1}, skew boxes]{ 3,2 | 2^2 }
  \end{example}

  There is no command for drawing shifted multitableau or multidiagrams.
  Such diagrams can be drawn using the
  \option{skew} key.  In the same way, it is possible to draw
  multitableaux and multidiagrams with components that are a mixture of
  non-skew, skew and shifted shapes.

  \begin{example}
    \Multitableau[skew={0,1|0,1}]{ 1236,45 | 7{10}{11},89} \par\bigskip
    \Multidiagram[skew={|0,1}]{ 3,2 | 2,1 }
  \end{example}

  \option[true]{tabloid}[true/false]

  Use the \option{tabloid} key for multitabloid diagrams and tableaux.

  \begin{example}
    \Multitableau[tabloid]{ 123,45 | 67,89} \par\bigskip
    \Multidiagram[tabloid]{ 4,2 | 2^2 }
  \end{example}

  \option[0]{xoffsets=offsets}[$|$-separated list of $x$-offsets]
  \option*[0]{yoffsets=offsets}[$|$-separated list of $y$-offsets]

  The \keyword{\Multitableau} and \keyword{\Multidiagram} commands try
  to draw the diagrams as compactly as possible with their first rows
  aligned, and the distance between the separators given by the value of
  \option{separation}. Use \option{xoffsets} to offset the
  $x$-coordinates of the component diagrams and \option{yoffsets} to
  offset their $y$-coordinates.

  \begin{example}
    \Multitableau[xoffsets={0.5,0.2}]{ 123,45 | 67,89} \par\bigskip
    \Multidiagram[yoffsets={0,1.5}]{ 4,2 | 2^2 }
  \end{example}

  When using \option{xoffsets} or \option{yoffsets}, you may want to use
  \option{no separators} to disable the separators between component
  shapes. For example, here is a homogeneous Garnir tableau using the
  style of~\cite{KMR:UniversalSpecht}, where no separators are used and
  the component tableaux are stacked vertically.

  \begin{example}[label=Ex:Garnir]
    \Multidiagram[entries=first, no separators,
      xoffsets={0|-1.1}, yoffsets={0|-1},
      styles={O={fill=orange!50}, P={fill=purple!50}},
      snob style={draw=red,double,thick,fill=white},
      snobs={|(O)24_{15}c_{14}, (P)26_{17}c_{16}, 27_{18},
        31_{11},  (P)33_{13}c_{12}}
    ]{1|7^2,4,1}
  \end{example}

  This is a slightly lazy way of drawing this multitableau because
  \option{entries=first} is used to automatically populate most of the
  entries and then \option{snobs=...} are used to rewrite the ``bricks''
  in the homogeneous Garnir belt.

  % --------------------------------------------------------------
  \section{Abacuses} \label{S:Abacus} \index{abacus} \indexcmd{Abacus}

  Gordon James introduced the abacus, with beads and runners, as another
  way to represent a partition~\cite{James:YoungD}. The abacus is
  particularly useful in modular representation theory, where the number
  of runners is the (quantum) characteristic. Abacus combinatorics are
  now used to help many algebras beyond their initial appearance in the
  representation theory of the symmetric groups.

  The partition corresponding to an abacus is given by counting the
  number of empty bead positions before each bead on the abacus. The
  \keyword{\Abacus} command, which draws an abacus representation of a
  partition, has a similar syntax to the tableau and diagram commands.

  \keyword{\Abacus (x,y) [options] {number of runners} {bead specifications}}

  As with other \aTableau commands, the $(x,y)$-coordinates are
  necessary if, and only if, the abacus is part of a
  \keyword{tikzpicture} environment.  The number of runners is a
  positive integer that specifies the number of abacus runners.

  % --------------------------------------------------------------
  \subsection{Bead specifications}

  In their simplest form, the bead specifications are a comma separated
  list of non-negative integers specifying a partition. Partitions can
  be given either as comma separated lists of non-negative integers, or
  using exponential notation. The full bead specifications add optional
  styling and bead labels.

  Here are some examples showing basic usage of the \keyword{\Abacus}
  command:

  \begin{Example}
  \aTabset{align=top}
  \Abacus{1}{2,1,1}     \qquad \Abacus{2}{2,1^2,0^3} \qquad \Abacus{3}{2,1^2,0^5} \qquad
  \Abacus{4}{2,1^2,0^5} \qquad \Abacus{5}{2,1^2,0^5} \qquad \Abacus{6}{2,1^2,0^5}
  \end{Example}

  The beads in an abacus correspond to the parts of the partition. Just
  as with the \keyword{\Tableau} command, each of the beads can be given
  optional styling, which can either be given by prepending a
  \keyword{*}, which applies the \option{abacus star style}, or by
  giving \TikZ styling-keys inside square brackets \keyword{[...]}.
  When exponential notation is used, the styles are applied to all of the
  associated beads.

  \begin{example}
  \aTabset{align=top}
  \Abacus{3}{3,*2^3,1}  \qquad
  \Abacus{4}{2,*1,[ball color=red]1^2, *1,
    [ball color=orange]1,[ball color=brown]0}
  \end{example}

  In particular, in the right-hand abacus, the two beads corresponding
  to \keyword{1^2} are coloured red, whereas the other three beads that
  correspond to a~$1$ are either orange or an off-white colour, which is
  the default \option{abacus star style}. As with the tableaux commands,
  we recommend using \keyword{\tikzset} to define more complicated bead styles.

  In addition to applying style to the individual beads on the abacus,
  \emph{labels} for any bead can (optionally) be given as a subscript in the
  bead specifications. Again, the same labels are applied to all of the
  corresponding beads when exponential notation is used.

  \begin{example}
  \Abacus{2}{5_5, 4_4, 3_3, 2_2,0_0} \qquad
  \Abacus[styles={R={ball color=red, text=yellow}}]
    {3}{*5^2_5, [R]4^3_4, 3, 2_2,*0^2_{oo}} \qquad
  \Abacus{4}{8^2_-,7_+,5_a,4^2_b,1_\times^3,1_d}
  \end{example}

  The bead labels are arbitrary, however, anything wider than the bead
  will spill into the abacus. Traditionally, abacus beads are not labelled,
  but you can give labels to as many beads as you like.

  You can remove beads that you have placed on the abacus using the \emph{no shade} style%
  \footnote{\TikZ \emph{does not} provide a way to disable
  a shading that is already in effect. The \emph{no shade} style, which is
  based on a \href{https://tex.stackexchange.com/a/85750/234252}{stack exchange post},
  is available whenever the \aTableau package is loaded.}.
  The point of being able to remove beads in this way is that this
  allows you to put something else in its place using the bead label.
  In such cases, you probably want to change the text colour, which is white
  by default.

  \begin{example}[label=Ex:NoShade]
  \Abacus[styles={C={no shade,text=red}}]
    {5}{8^2_-,7_+,5_a,4^2_b,[C]1_\bullet^3,1_d}
  \end{example}
  \index{no shade}

  As with tableau boxes, the bead labels are typeset in
  mathematics~mode by default. This can changed using
  \option{math entries} and \option{text entries}.
  \index{math entries!abacus}
  \index{text entries!abacus}

  % --------------------------------------------------------------
  \subsubsection{Abacus coordinates}\label{SS:abcoordinates}
  \index{xy@$(x,y)$}
  \index{Cartesian coordinates!abacus}

  The $(x,y)$-coordinates are required if, and only if, the
  \keyword{\Abacus}command is used inside a \keyword{tikzpicture}
  environment, in which case $(x,y)$ gives the coordinates of the
  ``outside corner'' of the box in row~$1$ and column~$1$ of the
  tableau. The following example shows how tableaux are placed using
  $(x,y)$-coordinates inside a \keyword{tikzpicture} environment. The
  example also shows that, by default, the beads and runners in an
  abacus are half a unit apart.

  \begin{example}
  \begin{tikzpicture}[add grid]
    \aTabset{abacus ends=_|}
    \Abacus(0,0.5)  [south]              {3}{2,1^2}
    \Abacus(0,2)    [east,  bead=red]    {3}{2,1^2}
    \Abacus(2.5,0.5)[north, bead=orange] {3}{2,1^2}
    \Abacus(3.5,-1) [west,  bead=brown]  {3}{2,1^2}
  \end{tikzpicture}
  \end{example}

  % --------------------------------------------------------------
  \subsection{Abacus keys}

  The abacus keys, or options, control the general appearance of the
  abacuses drawn by the \keyword{\Abacus} command. Many of the options,
  or keys, for tableaux and diagrams can be applied to abacuses, and we
  do not discuss all of these options here. For example, the
  \keyword{\Abacus} command accepts the following keys:
  \option{align}, \AddKeyLink{align-abacus}\AddKeyLink{halign-abacus}\AddKeyLink{valign-abacus}
  \option{math entries}, \AddKeyLink{math entries-abacus}
  \option{text entries}, \AddKeyLink{text entries-abacus}
  \option{name}, \AddKeyLink{name-abacus}
  \option{scale}, \AddKeyLink{scale-abacus}\AddKeyLink{xscale-abacus}\AddKeyLink{yscale-abacus}
  \option{styles}, \AddKeyLink{styles-abacus}
  \option{tikz after}, \AddKeyLink{tikz after-abacus}
  \option{tikz before}, and \AddKeyLink{tikz before-abacus}
  \option{tikzpicture}. \AddKeyLink{tikzpicture-abacus}
  See \autoref{S:keys} for the complete list.

  \option[-]{south}
  \option*[]{east}
  \option*[]{north}
  \option*[]{west}

  Just like tableaux, abacus have four different conventions.  By
  default, abacuses are drawn with the runners pointing to the south,
  from top to bottom, and runners increasing from left to right. Similar
  to tableaux, we number the rows $0,1,\dots$, and we number the
  runners, or columns, $0,1,\dots,e-1$, from left to right (when using
  \option{south}). If there are $e$ runners, then the bead in row~$r$
  and column~$c$ is in position $er+c$. See \autoref{Ex:BeadPositions}
  for an example.

  Using the keys \option{east}, \option{north} and \option{west},
  abacuses can be drawn so that the row index increases to the east,
  north, and west, respectively.

  \begin{Example}
\Abacus[south,entries=betas]{2}{4^2,2,1}\qquad\Abacus[east,entries=betas]{2}{4^2,2,1}\qquad
\Abacus[north,entries=betas]{2}{4^2,2,1}\qquad\Abacus[west,entries=betas]{2}{4^2,2,1}
  \end{Example}

  To make these conventions clearer, we have used
  \option{entries=betas}<abacus> to label the beads in these abacuses by
  their bead positions, or \emph{beta numbers}.

  \option[$-|$]{abacus ends}  [any pair of: $-$, \_\!\_, ., *, $|$, \textgreater]
  \option*[very thick, draw=\aTableauColour{Inner}]{abacus ends style=style}[\TikZ-styling]

  The \option{abacus ends} key sets the style used for the top and
  bottom of the abacus. The meaning of these six symbols is the
  following:

  \begin{quote}
    \begin{options}[labelwidth=1em, nosep]
      \item[*$-$] Adds a line that extends slightly past the first and last runners
      \item[*\_\!\_]  Adds a line between the first and last runners
      \item[*.]   Add dots to the end of each runner
      \item[**]   Add dots and an arrow to the end of each runner
      \item[*$|$] The runners do not have any additional decoration
      \item[*$>$] Add arrows to the end of each runner
    \end{options}
  \end{quote}

  Exactly two of these styling characters must be given when setting
  \option{abacus ends}. There are six options for the top of the abacus,
  and six options for the bottom, so there are 36 possible
  configurations (most of which will never be used...). We do not list
  all of them, but here are some examples.

  \begin{Example}
  \aTabset{align=top}
  \Abacus[abacus ends=-|]{3}{2,1^2,0^3}\qquad\Abacus[abacus ends=-_]{3}{2,1^2,0^3}\qquad
  \Abacus[abacus ends=_.]{3}{2,1^2,0^3}\qquad\Abacus[abacus ends=>>]{3}{2,1^2,0^3}\qquad
  \Abacus[abacus ends=>|]{3}{2,1^2,0^3}\qquad\Abacus[abacus ends=_-]{3}{2,1^2,0^3}\qquad
  \Abacus[abacus ends=>-]{3}{2,1^2,0^3}\qquad\Abacus[abacus ends=-*]{3}{2,1^2,0^3}
  \end{Example}

 As one more example, \emph{beta numbers} can be displayed on a single
 runner using:

  \begin{example}
  \Abacus[abacus ends=**, east, rows=7]{1}{2,1^2,*0^2}
  \end{example}

  The \option{abacus ends style} key changes the \TikZ-styling of the
  ends of the abacus. By default, this styling is the same as the
  styling for the abacus runners, which is set using
  \option{runner style}. Changes to \option{runner style} also change
  the style of the abacus ends, but their style is ultimately determined
  by \option{abacus ends style}.

  \begin{example}
  \Abacus[abacus ends style={aTableauMain,
      ultra thick}]{3}{2,1^2,0^3}\qquad
  \Abacus[runner style=cyan]{3}{2,1^2,0^3}\qquad
  \Abacus[abacus ends style=aTableauMain,
      runner style=cyan] {3}{2,1^2,0^3}
  \end{example}

  \option[ball color=\aTableauColour{StarStyle}, text=\aTableauColour{Main}]{abacus star style=style}
     [\TikZ-styling]

  The \option{abacus star style} key appends \TikZ-style keys to the
  abacus $*$-style:

  \begin{example}
  \Abacus{3}{5,*4_a,*1^2,0^4} \qquad
  \Abacus[abacus star style={ball color=yellow,text=black}]
    {3}{5,*4_a,*1^2,0^4}
  \end{example}

  \option[\aTableauColour{Main}]{bead=colour}[a \LaTeX colour]

  The \option{bead} key sets  the bead colour.

  \begin{example}
  \Abacus{2}{2^3,1} \qquad
  \Abacus[bead=black]{3}{4^3,2,1} \qquad
  \Abacus[bead=red]{4}{4^3,[ball color=blue]2,1}
  \end{example}

  The \option{bead} key changes the colour of every bead on the
  abacus. As shown, the colour of individual beads can be changed by
  applying the \TikZ-style \keyword{ball color} to the bead.

  \option[\textbackslash small]{bead font=font command}[\LaTeX\ font command]

  The \option{bead font} key sets the font used for the bead labels.
  By default, these labels are typeset as mathematics, so
  this is mainly useful for changing the font size (because, for example
  \keyword{\bfseries $1$} does not make the $1$ bold). It is only when
  you are using \option{text entries} that font commands like
  \keyword{\itshape} and \keyword{\bfseries} will have any effect.

  \begin{example}
  \Abacus{3}{3_3,2^2,1^3} \qquad
  \Abacus[bead font=\tiny]{3}{3_3,2^2,1^3} \qquad
  \Abacus[text entries, bead font=\bfseries]{3}{3_3,2^2,1^3}
  \end{example}

  The \option{scale} key can also be used to resize abacuses.

  \option[0.4]{bead size=number}[decimal number, distance in \unit{cm}]
  \option*[0.42]{bead sep=number}[decimal number, distance in \unit{cm}]
  \option*[0.42]{runner sep=number}[decimal number, distance in \unit{cm}]

  These three keys control the size of the beads and the distance
  between them, so they would normally be used together. By default, the
  \emph{diameter} of the abacus beads is \qty{0.4}{cm}, the distance
  between adjacent beads on the same runner is \qty{.42}{cm}, and the
  distance between consecutive runners is \qty{0.42}{cm}.

  \begin{example}
  \Abacus[bead size=0.6,
          bead sep=0.65,
          runner sep=0.65]
  {3}{3_3,2^2,1^3}
  \end{example}

  As with tableau, you can also change the size of the abacuses using
  \option{scale}<abacus>, \option{xscale}<abacus> and
  \option{yscale}<abacus>. Changing these keys globally affects both
  abacuses and tableaux.

  \option[]{bead style=style}[\TikZ-styling]

  The keys above are usually sufficient for fine-tuning the style of the
  abacus beads. Alternatively, the \option{bead style} key appends
  \TikZ-styling to the default bead style, which has many moving parts.

  \begin{example}
  \Abacus[bead style={shading=axis}, bead text=red] {3}{3_3,2^2,1^3} \qquad
  \Abacus[bead style={no shade, fill=aTableauMain}] {3}{3_3,2^2,1^3}
  \end{example}

  By default, the beads look like balls (they use the
  \TikZ-styling \keyword{shading=ball, ball color=}\aTableauColour{Main}).
  In contrast, the beads on the left look like go beads, and those
  on the right are flat.

  \option[white]{bead text=colour}[a \LaTeX~colour]

  The \option{bead text} key sets the text colour of the abacus labels.

  \begin{example}
  \Abacus[bead=blue, bead text=yellow] {3}{3_3,2^2,1^3}
  \end{example}

  \option[false]{beta numbers}[false/true]

  When \option{beta numbers} is set, the numbers in the bead
  specifications are entered as \emph{beta numbers}. The beta numbers
  give the bead positions, which are numbered $0,1,\dots$, starting from
  the first row and then continuing in this way in subsequent rows. More
  explicitly, if there are $e$ runners then the bead in row~$r$ and
  column~$c$ is in position $er+c$. Up to shift, bead numbers are first
  column hook lengths. In particular, beta numbers are pairwise
  distinct.

  \begin{example}
  \Abacus{3}{3,2^2} \qquad
  \Abacus[beta numbers]{3}{5,3,2} \qquad
  \Abacus[beta numbers]{3}{6,4,3,[ball color=red]0}
  \end{example}

  There is no inverse key to \option{beta numbers}, however,
  \option{beta numbers=false} will remove the requirement to use beta
  numbers in the bead specifications.

  \option[]{dotted cols=column indices}[list of runner indices]<abacus>
  \option*[]{dotted rows=row indices}[list of row indices]<abacus>

  The \option{dotted rows}<abacus> and \option{dotted cols}<abacus> keys
  draw abacuses where the specified rows and columns are replaced with
  dots. This makes it possible to draw \emph{generic} abacuses, where
  the number of rows and columns is not fully specified. Note that for
  abacuses, both the row and column indexing starts from~$0$.

  \begin{example}
  \Abacus[dotted rows={2,4}]{4}{8,5,4,3,2,1} \qquad
  \Abacus[dotted cols={1,3},abacus ends=-*]{6}{9,8,5,4,3^3,2}
  \end{example}

  As with tableau, consecutive rows and columns are treated together.

  \begin{example}
  \Abacus[dotted rows={1,2,3}]{4}{8,5,4,3,2,1} \qquad
  \Abacus[dotted cols={1,2,3},abacus ends=-*]{6}{9,8,5,4,3^3,2}
  \end{example}

  Comparing these two examples shows that the beads and runners in the
  dotted rows and columns are removed by \option{dotted rows}<abacus>
  and \option{dotted cols}<abacus>, respectively. Named coordinates,
  for use with \option{name}<abacus>, are not created for the bead and
  tick positions on the dotted rows and columns.

  The \option{dotted cols}<abacus> ad \option{dotted rows}<abacus> keys
  can be used together, in which case the \option{dotted rows}<abacus>
  are removed first, after which the \option{dotted cols}<abacus> are
  removed. The output may need to be adjusted if the dotted rows and
  columns overlap.  For example, consider the abacus:

  \begin{example}
  \Abacus[dotted cols={2,3},
          dotted rows={2,3}
  ]{6}{26,21^2,5^2,3,2^2,0^2}
  \end{example}

  The cross-hatched dots in the centre of the abacus are not ideal. The
  following trick addresses this:

  \begin{example}[lefthand width=0.16\textwidth, label=Ex:DottedAbacus]
  \Abacus[dotted cols={2,3}, dotted rows={2,3},
    tikz after = {
      \draw[aTableau/clearBoxes](A-1-1.south east)rectangle(A-4-4.north west);
      \draw[aTableau/dottedLine](A-1-1.south east)--(A-4-4.north west);
    }
  ]{6}{26,21^2,5^2,3,2^2,0^2}
  \end{example}

  \option[]{entries=value}[betas, residues, rows, shape]<abacus>

  Like the \option{entries} key for tableaux, the
  \option{entries}<abacus> key for abacuses provides a shortcut for
  labelling the abacus beads by some common choices. The possible
  choices are:

  \begin{itemize}

    \item\option{entries=betas}\index{entries!beta}\index{beta numbers!entries}
    \label{beta numbers!abacuses}
    Labels the beads by the corresponding \emph{beta numbers}.
  \begin{example}
  \Abacus[entries=betas]{3}{3,2^2,0^2} \qquad
  \Abacus[entries=betas]{4}{3,2^2,1^3,0^2}
  \end{example}

  \item\option{entries=residues}\index{entries!residues}
  Labels the beads by their residues. Unlike for tableaux, it is not
  necessary to specify the quantum characteristic \option{e}<abacus>,
  because this is determined by the number of runners and the
  \option{cartan}<abacus> type. You can override this by setting
  \option{e}<abacus> manually.
  \begin{example}
  \Abacus[entries=residues,]{3}{3,2^2,0^2} \qquad
  \Abacus[entries=residues, cartan=C]
    {4}{3,2^2,1^3,0^2}
  \end{example}

  \AddKeyLink{e-abacus}
  \AddKeyLink{charge-abacus}
  \AddKeyLink{cartan-abacus}
  \noindent
  As with (multi)tableau, use the \option{charge}<abacus> key to
  add an offset to the residues, and the \option{cartan}<abacus> key for
  residues of other Cartan types.

  \item\option{entries=rows}\index{entries!rows}
  Labels the beads by their rows index of the corresponding part in the
  partition.
  \begin{example}
  \Abacus[entries=rows]{3}{3,2^2,0^2} \qquad
  \Abacus[entries=rows]{4}{3,2^2,1^3,0^2}
  \end{example}
  \noindent
  As shown, the beads corresponding to zero parts of the corresponding
  partition are marked with a dash, since there is no corresponding row
  in partition for such beads.

  \item\option{entries=shape}\index{entries!partition}
  Labels the beads by their shape. That is, each bead is labelled by the
  corresponding by the corresponding part of the partition.
  \begin{example}
  \Abacus[entries=shape]{3}{3,2^2,0^2} \qquad
  \Abacus[entries=shape]{4}{3,2^2,1^3,0^2}
  \end{example}

  \end{itemize}

  \option[A]{name=text}{text}

  Just like the tableaux commands, the \keyword{abacus} command defines
  named coordinates for all of the \emph{beads} and \emph{ticks} on the
  abacus. By default, the named coordinates take the form
  $(\mathsf{A{-}r{-}c})$, for the bead or tick in row~$r$ and column~$c$ of the
  abacus. Note that for abacuses, both the row and column indexing starts
  from~$0$. The key \option{name} changes the prefix used for the names.
  For example, after \option{name=X} the bead and tick names take the
  form~$(\mathsf{X{-}r{-}c})$.

  \begin{Example}[label=Ex:BeadPositions]
    \Abacus[name=X,
      tikz after={
        \draw[orange,<-,thick] (X-0-0.west) -- ++(-1,0)
            node[blue,align=left,anchor=east] {Empty bead\\ position};
        \draw[red,<-,thick] (X-2-2.east) -- ++(2,0)
            node[blue,align=left,anchor=west] {First part of\\the partition};
    }]{3}{4^3,1^2}
  \end{Example}

  \option[]{no shade}[]

  The abacus balls are drawn with the \TikZ style
  \keyword{shading=ball}.  As noted in \autoref{Ex:NoShade}, you can
  disable the \TikZ-ball shading for individual beads using the \TikZ-style \keyword{no shade}. The
  \option{no shade} key disables the ball shading for \emph{all} beads,
  which allows you to fully customise how the abacus beads are
  displayed.

  \begin{example}
    \Abacus[no shade, bead text=red]{3}{3_1,2_2,2_3} \qquad
    \Abacus[bead style={no shade, shape=diamond,
      fill=brown, text=yellow}]{3}{3_1,2_2,2_3,0^2}
  \end{example}

  As the first example shows, if you use the \option{no shade} key then
  you need to define a new style for the abacus beads. As the second
  example shows, rather than using \option{no shade} as an
  \keyword{\Abacus} key, you can use it inside \option{bead style} as a
  \TikZ-style. (We include the \option{no shade} option partly for
  convenience and partly to advertise how to use it.)

  \option[]{rows=row index}[non-negative integer]<abacus>

  By default, the \keyword{\Abacus} command uses the smallest number of
  rows necessary to display all of the beads on the abacus. Use
  \option{rows}<abacus> to change the number of rows displayed. If you use
  the \option{rows}<abacus> key, then it is your responsibility to ensure that
  the abacus has enough rows to display all of the beads.

  \begin{example}
  \Abacus{3}{3,2^2,0^2} \qquad
  \Abacus[rows=4]{3}{3,2^2,0^2} \qquad
  \Abacus[rows=5]{3}{3,2^2,0^2}
  \end{example}

  \option[draw=\aTableauColour{Inner}]{runner=colour}[a \LaTeX colour]
  \option*[very thick]{runner style=style}[\TikZ-styling]

  Use \option{runner} to set the colour of abacus runners.

  \begin{example}
    \Abacus{3}{3,2^2} \qquad
    \Abacus[runner=red]{3}{5,3,2,0^2} \qquad
    \Abacus[runner=brown, bead=blue]
      {3}{6,4,3,[ball color=red]0^2}
  \end{example}

  The \option{runner style} key appends \TikZ-styling to the abacus
  runner style.

  \begin{example}
    \Abacus[runner style={dashed}]{3}{3,2^2} \quad
  \end{example}

  \option[]{runner labels=labels}{list of labels for the abacus runners}
  \option*[font=\textbackslash scriptsize, text=\aTableauColour{Inner}]{runner label style=style}[\TikZ-styling]

  The \option{runner labels} key adds labels to the each runner. An
  error if given if the number of labels does not match the number of
  runners. Use \option{runner label style} to change the style of the
  label.

  \begin{example}
  \Abacus[runner labels={1,2,0},
     entries=residues]{3}{3,2^2,0^2} \qquad
  \Abacus[runner labels={2,3,0,1}, entries=residues,
    runner label style={circle, draw=cyan, inner sep=0pt,
    minimum size=2.5mm}] {4}{5,3,2^2,0^2}
  \end{example}

  As with the \option{label} key for tableau, the runner labels are
  typeset as mathematics, by default.

  \option[\aTableauColour{Inner}]{tick=colour}[a \LaTeX colour]
  \option*[semithick]{tick style=style}[\TikZ-styling]
  \option*[0.1]{tick length=number}[length in \unit{cm}]

  These three keys control the ticks on the abacus runners, which mark
  the empty bead positions. Use \option{tick} to the colour of the
  ticks on the abacus runners. The other two key have the obvious meanings.

  \begin{example}
    \Abacus[tick=red]{3}{6,4,3,1,0} \qquad
    \Abacus[tick style={ultra thick,orange}]{3}{6,4,3,1,0} \qquad
    \Abacus[tick length=0.3]{3}{6,4,3,1,0} \qquad
  \end{example}

  % --------------------------------------------------------------
  \section{Examples from the literature}

  This section shows how to draw some tableaux that appear in the
  literature. All of these diagrams predate this package. The aim of
  this section is to show how these pictures can be drawn using this
  package.

  Diagrams like the following are common. Note the use of a ribbon to
  add the $\alpha$ in row~$6$ and column~$8$.

  \begin{example}
  \Diagram[french, no boxes, ribbons={*68_\alpha}]
    {15,13,12,9,8^2,5,3}
  \end{example}

  Mendes~ \cite{Mendes:hooks} uses tableaux like these to give a new
  proof of the Murnaghan-Nakayama rule.

  \begin{example}
    \tikzset{C/.style={circle,draw=blue,
              minimum size=3.5mm, fill=white,thick}}
    \Diagram[paths={
      [C]13_1cc_\bullet, [C]21_2r_\bullet,
      [C]23_3c_\bullet, [C]14_4
    }]{4,3,1} \qquad
    \Diagram[path style={rounded corners},
      paths={[C]12_1cr_\bullet, [C]14_2c_\bullet,
        [C]23_3c_\bullet, [C]31_4}]{4,3,1}
  \end{example}

  The following ribbon tableau comes from a nice paper by
  Fayers~\cite{Fayers:DyckTilings}, which is uses Dyck tilings to
  understand homogeneous Garnir relations for KLR algebras of type~$A$.

  \begin{example}
  \RibbonTableau[skew={1^2}, ukrainian, scale=0.7,
    ribbons={(draw=none)21_m}
  ]{
     12,16c_\circ c_\circ c_\circ r_\circ r_\circ *r,
     18crcccrr*r,19,1{10},1{11},1{12},
     22,29crcccrr*rccrcrr,
     2{12}ccrcrcccrr*r_lccrcrrr,
     31,*32,41,42,51,52,53,
     62cr,81,91,{10}1,{10}2,{11}1,{11}2
  }
  \end{example}

  The entries of a tableau are usually single characters, or numbers,
  but they can be more complicated. The paper
  \cite{bowman2023quiverpresentationsisomorphismshecke} contains
  tableaux with entries that are \emph{symbols}, which are produced by
  the \keyword{\Symb} command below.

  \begin{example}
  \newcommand\Symb[2]{\genfrac{}{}{0pt}{}{#1}{#2}}
  \Tableau[ukrainian, box font=\scriptsize,
     tikz after = {\draw[ultra thick, red]
        (A-2-3.east)--(A-2-2.south)--(A-3-2.west)
                    --(A-3-3.north)--cycle;
  }]{
     {\Symb11}{\Symb11}{\Symb11}{\Symb0s},
     {\Symb11}{\Symb11}{\Symb11}{\Symb0s},
     {\Symb11}{\Symb11}{\Symb{s^*}s}{\Symb0f},
     {\Symb11}{\Symb11}{\Symb11},
     {\Symb11}{\Symb11}
  }
  \end{example}

  The next example, also from
  \cite{bowman2023quiverpresentationsisomorphismshecke}, is a little
  more involved. Most of the diagram is drawn using \emph{pics} from
  \TikZ, which are a good way to add repeating features to a
  \keyword{tikzpicture}. In the code below, \emph{pics} are used to add
  the up and down strings to this drawing, taking as input the string
  colour, the starting box, the list of the boxes the strings goes
  through, and the final box.  At the risk of further obfuscation, this
  could be done a little more efficiently. From the perspective of this
  manual, the most interesting feature of this picture is that it is a
  non-trivial example that uses the named nodes for
  the tableaux boxes; see \autoref{E:NodeNames}.

\begin{Example}[label=Ex:waves]
\begin{tikzpicture}[
   Cap/.style={out=45, in=135}, % caps are drawn with to[Cap]
   Cup/.style={out=315,in=225}, % cups are drawn with to[Cup]
   every node/.style={font=\scriptsize},
   redribbon/.style={fill=red!40,opacity=0.4},
   % define pics for drawing the curvy up and down strings
   pics/ustring/.style n args = {4}{% {colour}{first coord}{wave coords}{last coord}
      code = {
          \def\last{#2} % should not be necessary
          \draw[thick,#1](T-|#2.north west)node[above=-3.5]{$\vee$}--(#2.north west)
              foreach [remember=\pt as \last (initially #2)] \pt in {#3}
                  {to[Cup](\last.north east) to[Cap] (\pt.north west)}
              to[Cup] (#4.north east) -- (#4.north east|-T)node[below=-3.5]{$\wedge$};
      }
   },
   pics/dstring/.style n args = {4}{% {colour}{first coord}{wave coords}{last coord}
      code = {
          \def\last{#2} % should not be necessary
          \draw[thick,#1](O-|#2.south west)node[below=-3.5]{$\wedge$}--(#2.south west)
              foreach [remember=\pt as \last (initially #2)] \pt in {#3}
                  {to[Cap](\last.south east) to[Cup] (\pt.south west)}
              to[Cap] (#4.south east) -- (#4.south east|-O)node[above=-3.5]{$\vee$};
      }
   },
 ]
 \coordinate (O) at (0,0); \coordinate (T) at (0,3.2);
 \draw(-2,0)--(2.5,0) (-2,3.2)--(2.5,3.2);
 \Diagram(0,0)[ukrainian, inner style=dotted, ribbons={(redribbon)24rcrcrc}]{5,4^3,3}
 \pic at (T) {ustring={gray}{A-1-5}{}{A-1-5}};  % up strings
 \pic at (T) {ustring={gray}{A-4-4}{}{A-4-4}};
 \pic at (T) {ustring={gray}{A-5-2}{A-4-3,A-3-4}{A-3-4}};
 \pic at (T) {ustring={gray}{A-5-3}{}{A-5-3}};
 \pic at (T) {ustring={red}{A-5-1}{A-4-2,A-3-3,A-2-4}{A-2-4}};
 \pic at (O) {dstring={gray}{A-1-1}{}{A-1-1}};  % down strings
 \pic at (O) {dstring={gray}{A-2-1}{A-1-2}{A-1-2}};
 \pic at (O) {dstring={gray}{A-3-1}{A-2-2,A-1-3}{A-1-3}};
 \pic at (O) {dstring={gray}{A-4-1}{A-3-2,A-2-3,A-1-4}{A-1-4}};
 \pic at (O) {dstring={gray}{A-5-1}{A-4-2,A-3-3,A-2-4,A-1-5}{A-1-5}};
\end{tikzpicture}
\end{Example}

  % --------------------------------------------------------------
  \section{Summary of \aTableau commands and options}\label{S:keys}
  \rowcolors{2}{LightSkyBlue!10}{}

  For this impatient, this section compactly lists the \aTableau
  commands and their options. The sections above show these commands and
  options are used, with examples.

  The \aTableau package provides the following commands:

  \begin{center}
    \begin{tabular}{Lll}\toprule
      \large\sffamily\color{MidnightBlue}Command
                     & \large\sffamily\color{MidnightBlue} Picture
                     & \large\sffamily\color{MidnightBlue} Section \\\midrule
      Abacus         & An abacus for a partition         & \autoref{S:Abacus}\\
      Diagram        & A Young diagram for a partition   & \autoref{S:Diagram}\\
      Multidiagram   & An $\ell$-tuple of Young diagrams & \autoref{S:Multitableau} \\
      Multitableau   & An  $\ell$-tuple of tableaux      & \autoref{S:Multitableau} \\
      RibbonTableau  & A ribbon tableau                  & \autoref{S:RibbonTableaux} \\
      ShiftedDiagram & A shifted Young diagram           & \autoref{S:Shifted} \\
      ShiftedTableau & A shifted tableau                 & \autoref{S:Shifted} \\
      SkewDiagram    & A skew Young diagram              & \autoref{S:Skew} \\
      SkewTableau    & A skew tableau                    & \autoref{S:Skew} \\
      Tableau        & A tableau                         & \autoref{S:Tableau} \\
      Tabloid        & A tabloid                         & \autoref{S:Tabloid} \\
      \bottomrule
    \end{tabular}
  \end{center}

  In addition, the \keyword{\aTabset} command is provided for setting
  can default \aTableau values for the \aTableau options.

  For easier reference, here is the list of the \aTableau keys, together
  with a quick explanation of what they do. The last two columns
  indicate whether the options apply to tableau (and diagrams), or to
  abacuses~--- or both!  Options starting with \keyword{(no)} are a pair
  of boolean options and their inverses. If not given a true or false
  value then they are implicitly set to true. Although not explicitly
  listed, spelling variations of keys from the American dialect
  of English are tolerated.

  Both the key names and the green ticks below are hyperlinks to the
  relevant sections of the manual.

  \newcommand\heading[1]{\textcolor{MidnightBlue}{\textsf{\large#1}}}
  \rowcolors{2}{LightSkyBlue!10}{}
  \begin{xltabular}{0.9\textwidth}{l>{\raggedright}Xcc}\toprule
    \heading{Key} & \heading{Meaning} & \heading{Tableau} & \heading{Abacus} \\
    \midrule
    \endhead
    \option{abacus ends}       & Set the top and bottom of the abacus                & \No  & \Yes\\
    \option{abacus ends style} & Set the style of the top and bottom of the abacus   & \No  & \Yes\\
    \option{align}             & Set baseline alignment of \aTableau diagram         & \Yes & \Yes[abacus]\\
    \option{australian}        & Use the Australian convention for tableaux          & \Yes & \No \\
    \option{bead}              & Set the colour of the abacus beads                  & \No  & \Yes \\
    \option{bead font}         & Set the font used for the abacus bead labels        & \No  & \Yes \\
    \option{bead sep}          & The distance between adjacent beads on a runner     & \No  & \Yes \\
    \option{bead size}         & The diameter of the abacus beads                    & \No  & \Yes \\
    \option{bead style}        & Set the style of the abacus beads                   & \No  & \Yes \\
    \option{bead text}         & Set the colour of the abacus bead labels            & \No  & \Yes \\
    \option{beta numbers}      & Enter bead specification as beta numbers            & \No  & \Yes \\
    \option*{border}           & Draw the tableau border                             & \Yes & \No \\
    \option{border colour}     & Set tableau border colour                           & \Yes & \No \\
    \option{border style}      & Set tableau border style                            & \Yes & \No \\
    \option{box height}        & Set tableau box height                              & \Yes & \No \\
    \option{box style}         & Set tableau box style                               & \Yes & \No \\
    \option{box width}         & Set tableau box width                               & \Yes & \No \\
    \option*{boxes}            & Draw inner walls around boxes                       & \Yes & \No \\
    \option{cartan}            & Set Cartan type for residues                        & \Yes & \Yes[abacus] \\
    \option{charge}            & Set charge for contents and residues                & \Yes & \Yes[abacus] \\
    \option{conjugate}         & Draw the conjugate tableau/diagram                  & \Yes & \No \\
    \option{e}                 & Set the quantum characteristic for residues         & \Yes & \Yes[abacus] \\
    \option{east}              & Make abacus grow in the easterly direction          & \Yes & \No \\
    \option{delimiters}        & Set delimiters for multitableau and multidiagrams   & \Yes & \No  \\
    \option{empty}             & Symbol for empty tableau in multitableau            & \Yes & \No \\
    \option{dotted cols}       & Specify columns to replace with dots                & \Yes & \Yes[abacus] \\
    \option{dotted rows}       & Specify rows to replace with dots                   & \Yes & \Yes[abacus] \\
    \option{english}           & Use the English convention for tableaux             & \Yes & \No \\
    \option{entries}           & Add entries to tableaux and abacuses                & \Yes & \Yes[abacus] \\
    \option{french}            & Use the French convention for tableaux              & \Yes & \No \\
    \option{halign}            & Horizontal alignment of box entries and bead labels & \Yes & \Yes[abacus] \\
    \option{inner style}       & Set style of inner tableau wall                     & \Yes & \No \\
    \option{inner wall}        & Set colour of inner tableau wall                    & \Yes & \No \\
    \option{label}             & Set tableau label                                   & \Yes & \No \\
    \option{label style}       & Set style of tableau labels                         & \Yes & \No \\
    \option{left delimiter}    & Set left delimiter for multitableau                 & \Yes & \No \\
    \option{math entries}      & Typeset box and label entries in math-mode          & \Yes & \Yes[abacus] \\
    \option{name}              & Prefix of named nodes                               & \Yes & \Yes[abacus] \\
    \option{north}             & Make abacus grow in the northerly direction         & \No  & \Yes\\
    \option{paths}             & Add paths to a diagram or tableau                   & \Yes & \No \\
    \option{path box}          & Adds default box entry for each path                & \Yes & \No \\
    \option{path box style}    & Sets style of boxes on each path                    & \Yes & \No \\
    \option{path style}        & Sets style of a path                                & \Yes & \No \\
    \option{ribbons}           & Add ribbons to tableau                              & \Yes & \No \\
    \option{ribbon box}        & Adds default box entry for each ribbon              & \Yes & \No \\
    \option{ribbon style}      & Set ribbon style                                    & \Yes & \No \\
    \option{ribbon box style}  & Sets style of boxes on each ribbon                  & \Yes & \No \\
    \option{right delimiter}   & List of labels for the abacus runners               & \No  & Yes \\
    \option{rows}              & Set number of rows in multitableau or abacus        & \Yes & \Yes[abacus] \\
    \option{runner}            & Set runner colour                                   & \No  & \Yes \\
    \option{runner labels}     & Set the colour of the abacus runners                & \No  & \Yes \\
    \option{runner label style}& Set the style for e abacus labels                   & \No  & \Yes \\
    \option{runner sep}        & Set distance between consecutive abacus runners     & \No  & \Yes \\
    \option{runner style}      & Set the colour of the abacus runners                & \No  & \Yes \\
    \option{scale}             & Sets the \aTableau scale                            & \Yes & \Yes[abacus] \\
    \option{script}            & Sets the \aTableau scale for subscripts             & \Yes & \Yes(scale-abacus) \\
    \option{scriptscript}      & Sets the \aTableau scale for subsubscripts          & \Yes & \Yes(scale-abacus) \\
    \option{separation}        & Set distance between components in multitableau     & \Yes & \No \\
    \option{separator colour}  & Set separator colour for multitableau               & \Yes & \No \\
    \option{separator}         & Set separator colour in multitableau                & \Yes & \No \\
    \option*{separators}       & Enable separators in multitableau                   & \Yes & \No \\
    \option{shifted}           & True for a shifted tableau                          & \Yes & \No \\
    \option{skew}              & Set inner skew shape                                & \Yes & \No \\
    \option*{skew border}      & Set skew border colour                              & \Yes & \No \\
    \option{skew border style} & Set \TikZ-style of the skew border                  & \Yes & \No \\
    \option*{skew boxes}       & Enable drawing of skew boxes                        & \Yes & \No \\
    \option{skew box style}    & Set \TikZ-style of the skew boxes                   & \Yes & \No \\
    \option{snobs}             & add snobs to tableau                                & \Yes & \No \\
    \option{snob box}          & Adds default box entry for each snob                & \Yes & \No \\
    \option{snob box style}    & Sets style of boxes on each snob                    & \Yes & \No \\
    \option{snob style}        & Sets style of a snob                                & \Yes & \No \\
    \option{south}             & Make abacus grow in the southerly direction         & \Yes & \No \\
    \option{star style}        & Set \TikZ-style of tableau $*$-nodes                & \Yes & \No \\
    \option{styles}            & Short-hand for defining single use \TikZ-styles     & \Yes & \Yes[abacus] \\
    \option{tabloid}           & True for a shifted tableau                          & \Yes & \No \\
    \option{text entries}      & Typeset box and label entries in text-mode          & \Yes & \Yes[abacus] \\
    \option{tick}              & Set the colour of the ticks on the abacus runner    & \No  & \Yes \\
    \option{tick length}       & Set the length of the ticks on the abacus runners   & \No  & \Yes \\
    \option{tick style}        & Set the style of the ticks on the abacus runners    & \No  & \Yes \\
    \option{tikz after}        & \TikZ-code injected after \aTableau picture         & \Yes & \Yes[abacus] \\
    \option{tikz before}       & \TikZ-code injected before \aTableau picture        & \Yes & \Yes[abacus] \\
    \option{tikzpicture}       & Sets tikzpicture environment keys                   & \Yes & \No \\
    \option{ukrainian}         & Use the Ukrainian convention for tableaux           & \Yes & \No \\
    \option{valign}            & Set vertical alignment                              & \Yes & \Yes[abacus] \\
    \option{west}              & Make abacus grow in the westerly direction          & \Yes & \No \\
    \option{xoffsets}          & Set $x$-offsets for components in a multitableau    & \Yes & \No \\
    \option{xscale}            & Set \aTableau $x$-scale                             & \Yes & \Yes(scale-abacus) \\
    \option{yoffsets}          & Set $y$-offsets for components in a multitableau    & \Yes & \No \\
    \option{yscale}            & Set \aTableau $y$-scale                             & \Yes & \Yes(scale-abacus) \\
    \bottomrule
  \end{xltabular}

  \emph{Most readers should ignore the following two subsections.}

  % --------------------------------------------------------------
  \subsection{\aTableau Colours}\label{S:Colours}

  The \aTableau package defines and uses the following colours:

  \rowcolors{2}{LightSkyBlue!10}{}
  \begin{center}
    \begin{tabular}{lll}\toprule
      \heading{Colour} & \heading{Name} & \heading{HTML}\\
       \midrule
       \aTableauColour{Inner}     & aTableauInner     & 0073E6 \\
       \aTableauColour{Main}      & aTableauMain      & 00008B \\
       \aTableauColour{SkewFill}  & aTableauSkewFill  & F8F8F8 \\
       \aTableauColour{Skew}      & aTableauSkew      & 818589 \\
       \aTableauColour{StarStyle} & aTableauStarStyle & E6F7FF \\
       \bottomrule
    \end{tabular}
  \end{center}

  All of these colours are defined as HTML colours using the
  \keyword{\definecolor} command provided by \ctan{xcolor}.

  We recommend using the key-value interface to change the
  colours used in the pictures created by \aTableau, rather than
  changing these colours. It is sometimes useful to be able to use these
  colour names directly in your own styles.

  % --------------------------------------------------------------
  \subsection{\aTableau styles}\label{S:TikZStyles}

  This package is little than a glorified interface to some \TikZ
  commands. Under the hood there are many custom \TikZ styles that control the
  pictures drawn by the \aTableau package. It is possible, and sometimes
  useful (see, for example, \autoref{Ex:DottedAbacus}), to use these
  styles in your own drawings.
  The key-value interface to the \aTableau commands is the recommended
  way of changing these styles. You can, if you want, change these
  styles directly using \keyword{\tikzset}, but there is a risk that you
  may break some of the \aTableau commands if you do this. If you do
  change these styles, we strongly recommend that you \emph{append} to
  these styles rather than setting them directly.

  \begin{center}
    \begin{xltabular}{0.9\textwidth}{ll}\toprule
       \heading{Names} & \heading{Adds \TikZ-styling to...}\\
       \midrule
       aTableau/innerWall        & Inner tableau walls\\
       aTableau/borderStyle      & Tableau borders \\
       aTableau/skewBorder       & Skew tableau borders \\
       aTableau/boxStyle         & Tableau boxes\\
       aTableau/skewBox          & Skew tableau boxes \\
       aTableau/pathBox          & Boxes on paths \\
       aTableau/ribbonBox        & Boxes on ribbons \\
       aTableau/snobBox          & Boxes on snobs \\
       aTableau/pathStyle        & Paths \\
       aTableau/ribbonStyle      & Ribbons \\
       aTableau/snobStyle        & Snobs \\
       aTableau/labelStyle       & Tableau labels \\
       aTableau/tableauStarStyle & Star style for tableaux \\
       aTableau/clearBoxes       & Erase material. Used by \option{dotted rows} and \option{dotted cols}\\
       aTableau/dottedLine       & Adds dotted lines. Used by \option{dotted rows} and \option{dotted cols}\\
       aTableau/separatorSymbol  & Separators in multitableau\\
       aTableau/separatorLine    & Separator lines in multitableau\\
       aTableau/leftDelimiter    & Left delimiter in multitableau\\
       aTableau/rightDelimiter   & Right delimiter in multitableau\\
       aTableau/beadStyle        & Abacus beads\\
       aTableau/runnerStyle      & Abacus runners\\
       aTableau/runnerLabelStyle & Style of abacus runner labels\\
       aTableau/abacusEnds       & Tops and bottoms of abacuses\\
       aTableau/abacusStarStyle  & Abacus stars \\
       aTableau/tickStyle        & Abacus ticks \\
       \bottomrule
    \end{xltabular}
  \end{center}

  Most of these styles are quite basic, but some of them are more
  involved, which is why we recommend not changing them directly. In
  some cases, such as the tableau boxes, the \TikZ styling does not
  fully control the rendering of these features. Most of these styles
  can be augmented by styles supplied by the key-value interface, which
  is another reason to avoid changing them directly.

  % --------------------------------------------------------------
  \section{Feature requests and bug reports}

  Please make any feature requests and bug reports on the package github
  page
  \begin{center}
    \href{https://github.com/AndrewMathas/aTableau/}{github.com/AndrewMathas/aTableau/}.
  \end{center}
  Please give as much detail as possible when making such requests.

  Bug reports must be a accompanied by a \textit{minimal working
  example}, which is the smallest possible amount of \LaTeX\ code that
  demonstrates the problem and compiles, unless you are reporting a
  problem that gives a compilation error in which case your example
  should give the error. This is necessary because if I cannot reproduce
  your problem, then it is unlikely that I will be able to fix it.


  % --------------------------------------------------------------
  \subsection*{Acknowledgements}

  Several features in the package have benefited from discussions on
  \url{tex.stackexchange}. In particular, I thank
    Max Chernoff,
    Lars Madsen,
    Clea F. Rees,
  and
    Jonathan P. Spratte.

% ----------------------------------------------------------------------
% bibliography
% \bibliography{papers}
% \bibliographystyle{andrew}

\begin{thebibliography}{10}

\bibitem{bowman2023quiverpresentationsisomorphismshecke}
\textsc{C.~Bowman, M.~D. Visscher, A.~Hazi, and C.~Stroppel},
  \href{https://arxiv.org/abs/2309.13695}{\textit{Quiver presentations and
  isomorphisms of Hecke categories and Khovanov arc algebras}}, 2023.
\newblock \href{http://arxiv.org/abs/2309.13695}{arXiv:2309.13695}.

\bibitem{DJM:cyc}
\textsc{R.~Dipper, G.~James, and A.~Mathas},
  \href{http://dx.doi.org/10.1007/PL00004665}{\textit{Cyclotomic {$q$}-{S}chur
  algebras}}, Math. Z., \textbf{229} (1998), 385--416.

\bibitem{EvseevMathas:DeformedKLR}
\textsc{A.~Evseev and A.~Mathas},
  \href{http://dx.doi.org/10.5802/art.8}{\textit{Content systems and
  deformations of cyclotomic {KLR} algebras of type~$A$ and~$C$}}, Annals of
  Representation Theory, \textbf{1} (2024), 193--297.
\newblock \href{http://arxiv.org/abs/2209.00134}{arXiv:2209.00134}.

\bibitem{Fayers:DyckTilings}
\textsc{M.~Fayers}, \textit{Dyck tilings and the homogeneous Garnir relations
  for graded Specht modules}, 2013, preprint.
\newblock \href{http://arxiv.org/abs/1309.6467}{arXiv:1309.6467}.

\bibitem{James:YoungD}
\textsc{G.~James}, \textit{Some combinatorial results involving {Y}oung
  diagrams}, Math. Proc. Cambridge Philos. Soc., \textbf{83} (1978), 1--10.

\bibitem{KMR:UniversalSpecht}
\textsc{A.~Kleshchev, A.~Mathas, and A.~Ram},
  \href{http://dx.doi.org/10.1112/plms/pds019}{\textit{Universal graded
  {S}pecht modules for cyclotomic {H}ecke algebras}}, Proc. Lond. Math. Soc.
  (3), \textbf{105} (2012), 1245--1289.
\newblock \href{http://arxiv.org/abs/1102.3519}{arXiv:1102.3519}.

\bibitem{Konvalinka:ShiftedHookLengths}
\textsc{M.~z. Konvalinka}, \href{http://dx.doi.org/10.37236/588}{\textit{The
  weighted hook length formula {III}: {S}hifted tableaux}}, Electron. J.
  Combin., \textbf{18} (2011), Paper 101, 29.

\bibitem{Macdonald}
\textsc{I.~G. Macdonald}, \textit{Symmetric functions and {H}all polynomials},
  Oxford Mathematical Monographs, The Clarendon Press Oxford University Press,
  New York, second~ed., 1995.
\newblock With contributions by A. Zelevinsky, Oxford Science Publications.

\bibitem{Mathas:ULect}
\textsc{A.~Mathas},
  \href{http://dx.doi.org/10.1090/ulect/015}{\textit{Iwahori-{H}ecke algebras
  and {S}chur algebras of the symmetric group}}, University Lecture Series,
  \textbf{15}, American Mathematical Society, Providence, RI, 1999.

\bibitem{Mendes:hooks}
\textsc{A.~Mendes}, \textit{The combinatorics of rim hook tableaux}, Australas.
  J. Combin., \textbf{73} (2019), 132--148.

\end{thebibliography}

% ----------------------------------------------------------------------
% indices
\printindex

\end{document}
